\section{未锚定账本事件的叙事补充：western-han}

\label{sec:anchor-batch-two-western-han}

以下事件均为此前正文尚未设置导航目标的账本记录。每一锚点置于来源、制度与地方条件的叙事中，避免把年表、政权分类或战报修辞当作社会全貌。

\subsection{制度与地方资源}

\eventanchor{WH-0139-ZHANG-QIAN}{前139年·张骞奉使西域寻求联络大月氏，途中为匈奴所留}账本记录汉与西域在前139年发生“张骞奉使西域寻求联络大月氏，途中为匈奴所留”。条目把背景概括为武帝希望在匈奴西方寻找可合作力量，以改变北方战略被动，并指出汉廷获得有关大宛、康居和大月氏的系统信息，西向外交与战争构想扩大。这个节点进入区域社会时，要经由文书、道路、军队、市场、寺院和家庭等不同中介；因此，政治行动的宣布、地方接收和日常生活的改变不能被视作同一时刻完成。

本条依据《史记·大宛列传》；《汉书·张骞李广利传》；《汉书·西域传上》，证据提示为“可靠纪年；出使月氏的直接战略目标有争议”。这些材料可支持年序、官方决定、战事或特定地点的线索，却难以直接统计所有人的财富、迁徙、信仰和动机。更稳妥的读法是区分诏令与实施、军报与人口、行政称谓与自我认同，并让地方材料的缺环限制解释的范围。

由此可见，该事件不是孤立的年表点，而是资源、信息与权力关系重新配置的一环。不同地区的官员、社群和家庭可能以不同方式承担征发、保护、迁移或宗教活动；现存来源只能照亮其中一部分，不能替沉默者制造统一的经验。

\eventanchor{WH-0136-FIVE-CLASSICS}{前136年·武帝置五经博士，国家经学官职获得制度位置}账本记录汉在前136年发生“武帝置五经博士，国家经学官职获得制度位置”。条目把背景概括为宫廷礼制与政治合法性需要可训练、可任用的经典解释者，并指出经师进入官僚制度，但黄老、方术和法令传统并未因此消失。这个节点进入区域社会时，要经由文书、道路、军队、市场、寺院和家庭等不同中介；因此，政治行动的宣布、地方接收和日常生活的改变不能被视作同一时刻完成。

本条依据《汉书·武帝纪》；《汉书·儒林传》；《汉书·百官公卿表》，证据提示为“可靠纪年；学官规模与课程有争议”。这些材料可支持年序、官方决定、战事或特定地点的线索，却难以直接统计所有人的财富、迁徙、信仰和动机。更稳妥的读法是区分诏令与实施、军报与人口、行政称谓与自我认同，并让地方材料的缺环限制解释的范围。

由此可见，该事件不是孤立的年表点，而是资源、信息与权力关系重新配置的一环。不同地区的官员、社群和家庭可能以不同方式承担征发、保护、迁移或宗教活动；现存来源只能照亮其中一部分，不能替沉默者制造统一的经验。

\eventanchor{WH-0135-MINYUE}{前135年·闽越攻南越，汉廷遣军干预并调整东南封国关系}账本记录汉与闽越、南越在前135年发生“闽越攻南越，汉廷遣军干预并调整东南封国关系”。条目把背景概括为闽越与南越争夺东南资源和海陆通道，南越向汉求援，并指出汉廷以援军和封王调换介入东南，扩张前的羁縻秩序开始改变。这个节点进入区域社会时，要经由文书、道路、军队、市场、寺院和家庭等不同中介；因此，政治行动的宣布、地方接收和日常生活的改变不能被视作同一时刻完成。

本条依据《汉书·武帝纪》；《史记·南越列传》；《汉书·闽粤传》，证据提示为“可靠纪年；出兵规模有争议”。这些材料可支持年序、官方决定、战事或特定地点的线索，却难以直接统计所有人的财富、迁徙、信仰和动机。更稳妥的读法是区分诏令与实施、军报与人口、行政称谓与自我认同，并让地方材料的缺环限制解释的范围。

由此可见，该事件不是孤立的年表点，而是资源、信息与权力关系重新配置的一环。不同地区的官员、社群和家庭可能以不同方式承担征发、保护、迁移或宗教活动；现存来源只能照亮其中一部分，不能替沉默者制造统一的经验。

\eventanchor{WH-0134-XIAOLIAN}{前134年·武帝令郡国举孝廉，察举成为选官的重要渠道}账本记录汉在前134年发生“武帝令郡国举孝廉，察举成为选官的重要渠道”。条目把背景概括为中央扩张官僚体系，需要从郡国取得兼具名望与行政潜能的人才，并指出地方评价进入中央选官程序，士人德行语言与官僚流动逐渐相连。这个节点进入区域社会时，要经由文书、道路、军队、市场、寺院和家庭等不同中介；因此，政治行动的宣布、地方接收和日常生活的改变不能被视作同一时刻完成。

本条依据《汉书·武帝纪》；《汉书·儒林传》；《后汉书·百官志》，证据提示为“可靠纪年；早期执行频率有争议”。这些材料可支持年序、官方决定、战事或特定地点的线索，却难以直接统计所有人的财富、迁徙、信仰和动机。更稳妥的读法是区分诏令与实施、军报与人口、行政称谓与自我认同，并让地方材料的缺环限制解释的范围。

由此可见，该事件不是孤立的年表点，而是资源、信息与权力关系重新配置的一环。不同地区的官员、社群和家庭可能以不同方式承担征发、保护、迁移或宗教活动；现存来源只能照亮其中一部分，不能替沉默者制造统一的经验。

\eventanchor{WH-0133-MAYI}{前133年·马邑伏击计划泄露，汉匈关系转入长期大战}账本记录汉与匈奴在前133年发生“马邑伏击计划泄露，汉匈关系转入长期大战”。条目把背景概括为和亲互市未能解决边境冲突，武帝朝形成以军事主动权制约匈奴的共识，并指出双方互市关系受损，汉廷开始常态化组织大规模骑兵远征。这个节点进入区域社会时，要经由文书、道路、军队、市场、寺院和家庭等不同中介；因此，政治行动的宣布、地方接收和日常生活的改变不能被视作同一时刻完成。

本条依据《史记·匈奴列传》；《汉书·武帝纪》；《汉书·匈奴传》，证据提示为“可靠纪年；诱敌方案细节有异文”。这些材料可支持年序、官方决定、战事或特定地点的线索，却难以直接统计所有人的财富、迁徙、信仰和动机。更稳妥的读法是区分诏令与实施、军报与人口、行政称谓与自我认同，并让地方材料的缺环限制解释的范围。

由此可见，该事件不是孤立的年表点，而是资源、信息与权力关系重新配置的一环。不同地区的官员、社群和家庭可能以不同方式承担征发、保护、迁移或宗教活动；现存来源只能照亮其中一部分，不能替沉默者制造统一的经验。

\eventanchor{WH-0129-WEI-QING}{前129年·卫青等四将出塞，龙城之战取得武帝朝首场显著战果}账本记录汉与匈奴在前129年发生“卫青等四将出塞，龙城之战取得武帝朝首场显著战果”。条目把背景概括为马邑失败后汉廷试探性多路出击，需验证骑兵远征能力，并指出卫青崛起，汉军由被动守边逐渐转向主动袭击匈奴部众。这个节点进入区域社会时，要经由文书、道路、军队、市场、寺院和家庭等不同中介；因此，政治行动的宣布、地方接收和日常生活的改变不能被视作同一时刻完成。

本条依据《史记·卫将军骠骑列传》；《汉书·卫青霍去病传》；《资治通鉴》卷十八，证据提示为“核心事实可靠；四路行军和斩获数有争议”。这些材料可支持年序、官方决定、战事或特定地点的线索，却难以直接统计所有人的财富、迁徙、信仰和动机。更稳妥的读法是区分诏令与实施、军报与人口、行政称谓与自我认同，并让地方材料的缺环限制解释的范围。

由此可见，该事件不是孤立的年表点，而是资源、信息与权力关系重新配置的一环。不同地区的官员、社群和家庭可能以不同方式承担征发、保护、迁移或宗教活动；现存来源只能照亮其中一部分，不能替沉默者制造统一的经验。

\eventanchor{WH-0127-SHUOFANG}{前127年·卫青收复河南地，汉置朔方郡并修筑防御体系}账本记录汉与匈奴在前127年发生“卫青收复河南地，汉置朔方郡并修筑防御体系”。条目把背景概括为匈奴据河套可威胁关中北缘，汉廷须建立前推据点和补给基地，并指出河套纳入汉郡县与屯戍网络，边防成本和徙民压力同步增加。这个节点进入区域社会时，要经由文书、道路、军队、市场、寺院和家庭等不同中介；因此，政治行动的宣布、地方接收和日常生活的改变不能被视作同一时刻完成。

本条依据《史记·卫将军骠骑列传》；《汉书·武帝纪》；《汉书·地理志》，证据提示为“可靠纪年；郡界和筑塞范围有争议”。这些材料可支持年序、官方决定、战事或特定地点的线索，却难以直接统计所有人的财富、迁徙、信仰和动机。更稳妥的读法是区分诏令与实施、军报与人口、行政称谓与自我认同，并让地方材料的缺环限制解释的范围。

由此可见，该事件不是孤立的年表点，而是资源、信息与权力关系重新配置的一环。不同地区的官员、社群和家庭可能以不同方式承担征发、保护、迁移或宗教活动；现存来源只能照亮其中一部分，不能替沉默者制造统一的经验。

\eventanchor{WH-0124-MO-NAN}{前124年·卫青夜袭右贤王，扩大漠南作战成果}账本记录汉与匈奴在前124年发生“卫青夜袭右贤王，扩大漠南作战成果”。条目把背景概括为朔方前推后汉军获得更近的出塞基地，武帝要求持续压迫匈奴右翼，并指出汉军提高草原作战经验，但远征耗费加重对马匹、粮秣和财政的需求。这个节点进入区域社会时，要经由文书、道路、军队、市场、寺院和家庭等不同中介；因此，政治行动的宣布、地方接收和日常生活的改变不能被视作同一时刻完成。

本条依据《史记·卫将军骠骑列传》；《汉书·卫青霍去病传》；《汉书·匈奴传》，证据提示为“可靠纪年；俘获和兵数有争议”。这些材料可支持年序、官方决定、战事或特定地点的线索，却难以直接统计所有人的财富、迁徙、信仰和动机。更稳妥的读法是区分诏令与实施、军报与人口、行政称谓与自我认同，并让地方材料的缺环限制解释的范围。

由此可见，该事件不是孤立的年表点，而是资源、信息与权力关系重新配置的一环。不同地区的官员、社群和家庭可能以不同方式承担征发、保护、迁移或宗教活动；现存来源只能照亮其中一部分，不能替沉默者制造统一的经验。

\eventanchor{WH-0121-HEXI}{前121年春夏·霍去病两次出陇西、河西，击破浑邪、休屠部}账本记录汉与匈奴在前121年春夏发生“霍去病两次出陇西、河西，击破浑邪、休屠部”。条目把背景概括为汉廷试图切断匈奴与西羌、西域的联络，霍去病率轻骑深入河西，并指出河西走廊军事格局转变，汉可进一步设置郡县并经营西向通道。这个节点进入区域社会时，要经由文书、道路、军队、市场、寺院和家庭等不同中介；因此，政治行动的宣布、地方接收和日常生活的改变不能被视作同一时刻完成。

本条依据《史记·卫将军骠骑列传》；《汉书·霍去病传》；《汉书·匈奴传》，证据提示为“可靠纪年；战果与行军距离有争议”。这些材料可支持年序、官方决定、战事或特定地点的线索，却难以直接统计所有人的财富、迁徙、信仰和动机。更稳妥的读法是区分诏令与实施、军报与人口、行政称谓与自我认同，并让地方材料的缺环限制解释的范围。

由此可见，该事件不是孤立的年表点，而是资源、信息与权力关系重新配置的一环。不同地区的官员、社群和家庭可能以不同方式承担征发、保护、迁移或宗教活动；现存来源只能照亮其中一部分，不能替沉默者制造统一的经验。

\eventanchor{WH-0121-HUNYE-SURRENDER}{前121年秋·浑邪王率部众降汉，汉廷安置降众于边郡}账本记录汉与匈奴在前121年秋发生“浑邪王率部众降汉，汉廷安置降众于边郡”。条目把背景概括为河西连续失利使浑邪王与匈奴单于关系恶化，汉军提供可选择的保护与赏赐，并指出汉获得河西控制的政治支点，降人安置也带来粮食、身份与边防管理问题。这个节点进入区域社会时，要经由文书、道路、军队、市场、寺院和家庭等不同中介；因此，政治行动的宣布、地方接收和日常生活的改变不能被视作同一时刻完成。

本条依据《史记·卫将军骠骑列传》；《汉书·霍去病传》；《汉书·地理志》，证据提示为“核心事实可靠；降众数量有争议”。这些材料可支持年序、官方决定、战事或特定地点的线索，却难以直接统计所有人的财富、迁徙、信仰和动机。更稳妥的读法是区分诏令与实施、军报与人口、行政称谓与自我认同，并让地方材料的缺环限制解释的范围。

由此可见，该事件不是孤立的年表点，而是资源、信息与权力关系重新配置的一环。不同地区的官员、社群和家庭可能以不同方式承担征发、保护、迁移或宗教活动；现存来源只能照亮其中一部分，不能替沉默者制造统一的经验。

\eventanchor{WH-0120-HEXI-COMMANDERIES}{前120年至前111年·河西陆续设置武威、酒泉、张掖、敦煌等郡，经营走廊}账本记录汉在前120年至前111年发生“河西陆续设置武威、酒泉、张掖、敦煌等郡，经营走廊”。条目把背景概括为河西战胜需要转化为持续行政和交通控制，中央须安置军民并保障道路，并指出走廊成为汉与西域往来的制度化通道，但郡县实际覆盖呈现边塞与绿洲差异。这个节点进入区域社会时，要经由文书、道路、军队、市场、寺院和家庭等不同中介；因此，政治行动的宣布、地方接收和日常生活的改变不能被视作同一时刻完成。

本条依据《汉书·地理志》；《汉书·武帝纪》；居延汉简（河西边郡文书），证据提示为“设置次序可靠；各郡建置年份有争议”。这些材料可支持年序、官方决定、战事或特定地点的线索，却难以直接统计所有人的财富、迁徙、信仰和动机。更稳妥的读法是区分诏令与实施、军报与人口、行政称谓与自我认同，并让地方材料的缺环限制解释的范围。

由此可见，该事件不是孤立的年表点，而是资源、信息与权力关系重新配置的一环。不同地区的官员、社群和家庭可能以不同方式承担征发、保护、迁移或宗教活动；现存来源只能照亮其中一部分，不能替沉默者制造统一的经验。

\eventanchor{WH-0119-MOBEI}{前119年·卫青、霍去病北出漠北，匈奴单于远遁}账本记录汉与匈奴在前119年发生“卫青、霍去病北出漠北，匈奴单于远遁”。条目把背景概括为汉已取得河套和河西基地，试图以大规模会战削弱单于本部，并指出匈奴未被消灭但北退，汉军马匹和财政损耗巨大，战争政策引出财政集权。这个节点进入区域社会时，要经由文书、道路、军队、市场、寺院和家庭等不同中介；因此，政治行动的宣布、地方接收和日常生活的改变不能被视作同一时刻完成。

本条依据《史记·卫将军骠骑列传》；《汉书·武帝纪》；《汉书·匈奴传》，证据提示为“可靠纪年；双方损失与战果有争议”。这些材料可支持年序、官方决定、战事或特定地点的线索，却难以直接统计所有人的财富、迁徙、信仰和动机。更稳妥的读法是区分诏令与实施、军报与人口、行政称谓与自我认同，并让地方材料的缺环限制解释的范围。

由此可见，该事件不是孤立的年表点，而是资源、信息与权力关系重新配置的一环。不同地区的官员、社群和家庭可能以不同方式承担征发、保护、迁移或宗教活动；现存来源只能照亮其中一部分，不能替沉默者制造统一的经验。

\eventanchor{WH-0119-SALT-IRON-COIN}{前119年·武帝朝收铸五铢钱并强化盐铁专营，由桑弘羊等主持财政措施}账本记录汉在前119年发生“武帝朝收铸五铢钱并强化盐铁专营，由桑弘羊等主持财政措施”。条目把背景概括为漠北远征和边郡建设耗费巨大，私铸与地方资源控制限制中央财力，并指出国家直接介入盐铁和铸币，形成持续争论的财政官营体系。这个节点进入区域社会时，要经由文书、道路、军队、市场、寺院和家庭等不同中介；因此，政治行动的宣布、地方接收和日常生活的改变不能被视作同一时刻完成。

本条依据《史记·平准书》；《汉书·食货志》；西汉五铢钱范与钱币窖藏，证据提示为“年代大体可靠；措施分期和执行范围有争议”。这些材料可支持年序、官方决定、战事或特定地点的线索，却难以直接统计所有人的财富、迁徙、信仰和动机。更稳妥的读法是区分诏令与实施、军报与人口、行政称谓与自我认同，并让地方材料的缺环限制解释的范围。

由此可见，该事件不是孤立的年表点，而是资源、信息与权力关系重新配置的一环。不同地区的官员、社群和家庭可能以不同方式承担征发、保护、迁移或宗教活动；现存来源只能照亮其中一部分，不能替沉默者制造统一的经验。

\eventanchor{WH-0118-UNIFORM-TRANSPORT}{前118年前后·桑弘羊推行均输、平准以调运物资和调节价格}账本记录汉在前118年前后发生“桑弘羊推行均输、平准以调运物资和调节价格”。条目把背景概括为远征与都城供应需要将各地贡赋转化为可调度的物资，运输差价被商人利用，并指出中央财政取得更多流通信息和收益，也扩大官商竞争及地方执行成本。这个节点进入区域社会时，要经由文书、道路、军队、市场、寺院和家庭等不同中介；因此，政治行动的宣布、地方接收和日常生活的改变不能被视作同一时刻完成。

本条依据《史记·平准书》；《汉书·食货志》；《盐铁论·本议》，证据提示为“制度存在可靠；实施地域与效果有争议”。这些材料可支持年序、官方决定、战事或特定地点的线索，却难以直接统计所有人的财富、迁徙、信仰和动机。更稳妥的读法是区分诏令与实施、军报与人口、行政称谓与自我认同，并让地方材料的缺环限制解释的范围。

由此可见，该事件不是孤立的年表点，而是资源、信息与权力关系重新配置的一环。不同地区的官员、社群和家庭可能以不同方式承担征发、保护、迁移或宗教活动；现存来源只能照亮其中一部分，不能替沉默者制造统一的经验。

\eventanchor{WH-0115-ZHANG-QIAN-SECOND}{前115年·张骞再使乌孙并派副使赴大宛、康居等地}账本记录汉与西域诸国在前115年发生“张骞再使乌孙并派副使赴大宛、康居等地”。条目把背景概括为河西走廊初步可通，汉廷希望以乌孙牵制匈奴并建立多边联系，并指出西域诸国进入汉廷较稳定的外交视野，贸易、贡使和军事信息往来增加。这个节点进入区域社会时，要经由文书、道路、军队、市场、寺院和家庭等不同中介；因此，政治行动的宣布、地方接收和日常生活的改变不能被视作同一时刻完成。

本条依据《史记·大宛列传》；《汉书·张骞李广利传》；《汉书·西域传上·乌孙国》，证据提示为“可靠纪年；使团路线与回报材料有争议”。这些材料可支持年序、官方决定、战事或特定地点的线索，却难以直接统计所有人的财富、迁徙、信仰和动机。更稳妥的读法是区分诏令与实施、军报与人口、行政称谓与自我认同，并让地方材料的缺环限制解释的范围。

由此可见，该事件不是孤立的年表点，而是资源、信息与权力关系重新配置的一环。不同地区的官员、社群和家庭可能以不同方式承担征发、保护、迁移或宗教活动；现存来源只能照亮其中一部分，不能替沉默者制造统一的经验。

\eventanchor{WH-0112-NANYUE-CAMPAIGN}{前112年·南越相吕嘉杀汉使，武帝发兵征南越}账本记录汉与南越在前112年发生“南越相吕嘉杀汉使，武帝发兵征南越”。条目把背景概括为南越王朝内亲汉与反汉力量冲突，汉廷试图将藩属关系改造为直接影响，并指出汉军大规模南征开始，岭南将由朝贡关系转向郡县化。这个节点进入区域社会时，要经由文书、道路、军队、市场、寺院和家庭等不同中介；因此，政治行动的宣布、地方接收和日常生活的改变不能被视作同一时刻完成。

本条依据《史记·南越列传》；《汉书·武帝纪》；《汉书·南粤传》，证据提示为“可靠纪年；南越内部派系与死者细节有争议”。这些材料可支持年序、官方决定、战事或特定地点的线索，却难以直接统计所有人的财富、迁徙、信仰和动机。更稳妥的读法是区分诏令与实施、军报与人口、行政称谓与自我认同，并让地方材料的缺环限制解释的范围。

由此可见，该事件不是孤立的年表点，而是资源、信息与权力关系重新配置的一环。不同地区的官员、社群和家庭可能以不同方式承担征发、保护、迁移或宗教活动；现存来源只能照亮其中一部分，不能替沉默者制造统一的经验。

\eventanchor{WH-0111-NANYUE-ANNEXATION}{前111年·汉军入番禺，灭南越并置南海等郡}账本记录汉与南越在前111年发生“汉军入番禺，灭南越并置南海等郡”。条目把背景概括为吕嘉政变使汉廷拥有军事干预的直接理由，水陆军自多路进入岭南，并指出岭南纳入汉朝行政版图；郡县、旧越首领和移民群体的互动需按地域分辨。这个节点进入区域社会时，要经由文书、道路、军队、市场、寺院和家庭等不同中介；因此，政治行动的宣布、地方接收和日常生活的改变不能被视作同一时刻完成。

本条依据《史记·南越列传》；《汉书·武帝纪》；《汉书·地理志》，证据提示为“可靠纪年；郡数及地方控制速度有争议”。这些材料可支持年序、官方决定、战事或特定地点的线索，却难以直接统计所有人的财富、迁徙、信仰和动机。更稳妥的读法是区分诏令与实施、军报与人口、行政称谓与自我认同，并让地方材料的缺环限制解释的范围。

由此可见，该事件不是孤立的年表点，而是资源、信息与权力关系重新配置的一环。不同地区的官员、社群和家庭可能以不同方式承担征发、保护、迁移或宗教活动；现存来源只能照亮其中一部分，不能替沉默者制造统一的经验。

\eventanchor{WH-0110-FENGSHAN}{前110年·武帝东巡泰山举行封禅，重塑皇帝祭祀权威}账本记录汉在前110年发生“武帝东巡泰山举行封禅，重塑皇帝祭祀权威”。条目把背景概括为北方和南方军事扩张增强皇帝自我表述的资源，方士与儒生对礼制解释竞争，并指出封禅将战功、天命和祭祀相连，郊祀制度成为中央政治与财政投入的议题。这个节点进入区域社会时，要经由文书、道路、军队、市场、寺院和家庭等不同中介；因此，政治行动的宣布、地方接收和日常生活的改变不能被视作同一时刻完成。

本条依据《史记·封禅书》；《汉书·郊祀志》；《汉书·武帝纪》，证据提示为“可靠纪年；礼仪复原有争议”。这些材料可支持年序、官方决定、战事或特定地点的线索，却难以直接统计所有人的财富、迁徙、信仰和动机。更稳妥的读法是区分诏令与实施、军报与人口、行政称谓与自我认同，并让地方材料的缺环限制解释的范围。

由此可见，该事件不是孤立的年表点，而是资源、信息与权力关系重新配置的一环。不同地区的官员、社群和家庭可能以不同方式承担征发、保护、迁移或宗教活动；现存来源只能照亮其中一部分，不能替沉默者制造统一的经验。

\eventanchor{WH-0109-DIAN}{前109年·汉军入滇，滇王降，设益州郡}账本记录汉与滇在前109年发生“汉军入滇，滇王降，设益州郡”。条目把背景概括为汉廷欲控制西南道路和资源，西南夷政权在周边战事压力下调整关系，并指出益州郡成为西南行政名目，滇池地区王权与汉制并存的具体形态仍须考古检验。这个节点进入区域社会时，要经由文书、道路、军队、市场、寺院和家庭等不同中介；因此，政治行动的宣布、地方接收和日常生活的改变不能被视作同一时刻完成。

本条依据《史记·西南夷列传》；《汉书·武帝纪》；《汉书·西南夷两粤朝鲜传》，证据提示为“可靠纪年；军事路线与地方统治程度有争议”。这些材料可支持年序、官方决定、战事或特定地点的线索，却难以直接统计所有人的财富、迁徙、信仰和动机。更稳妥的读法是区分诏令与实施、军报与人口、行政称谓与自我认同，并让地方材料的缺环限制解释的范围。

由此可见，该事件不是孤立的年表点，而是资源、信息与权力关系重新配置的一环。不同地区的官员、社群和家庭可能以不同方式承担征发、保护、迁移或宗教活动；现存来源只能照亮其中一部分，不能替沉默者制造统一的经验。

\eventanchor{WH-0108-WIMAN-CHOSEON}{前108年·汉灭卫满朝鲜，设乐浪等郡}账本记录汉与卫满朝鲜在前108年发生“汉灭卫满朝鲜，设乐浪等郡”。条目把背景概括为朝鲜政权与汉使、周边部族的关系恶化，汉廷试图控制东北交通和外交，并指出乐浪等郡成为汉在半岛北部的行政据点，其实际辖境和文化影响不可等同于全半岛统属。这个节点进入区域社会时，要经由文书、道路、军队、市场、寺院和家庭等不同中介；因此，政治行动的宣布、地方接收和日常生活的改变不能被视作同一时刻完成。

本条依据《史记·朝鲜列传》；《汉书·武帝纪》；《汉书·地理志》，证据提示为“核心事件可靠；郡数、位置和存续时间有争议”。这些材料可支持年序、官方决定、战事或特定地点的线索，却难以直接统计所有人的财富、迁徙、信仰和动机。更稳妥的读法是区分诏令与实施、军报与人口、行政称谓与自我认同，并让地方材料的缺环限制解释的范围。

由此可见，该事件不是孤立的年表点，而是资源、信息与权力关系重新配置的一环。不同地区的官员、社群和家庭可能以不同方式承担征发、保护、迁移或宗教活动；现存来源只能照亮其中一部分，不能替沉默者制造统一的经验。

\eventanchor{WH-0104-TAICHU-CALENDAR}{前104年·武帝改用太初历，调整岁首与历法计算}账本记录汉在前104年发生“武帝改用太初历，调整岁首与历法计算”。条目把背景概括为旧历与天象、祭祀和行政岁首的对应引发争论，中央需要统一的历日权威，并指出太初历为官府纪年和祭祀提供新框架，后世对其计算细节仍有不同复原。这个节点进入区域社会时，要经由文书、道路、军队、市场、寺院和家庭等不同中介；因此，政治行动的宣布、地方接收和日常生活的改变不能被视作同一时刻完成。

本条依据《汉书·武帝纪》；《汉书·律历志》；《史记·历书》，证据提示为“可靠纪年；具体算法与历法传播有争议”。这些材料可支持年序、官方决定、战事或特定地点的线索，却难以直接统计所有人的财富、迁徙、信仰和动机。更稳妥的读法是区分诏令与实施、军报与人口、行政称谓与自我认同，并让地方材料的缺环限制解释的范围。

由此可见，该事件不是孤立的年表点，而是资源、信息与权力关系重新配置的一环。不同地区的官员、社群和家庭可能以不同方式承担征发、保护、迁移或宗教活动；现存来源只能照亮其中一部分，不能替沉默者制造统一的经验。

\eventanchor{WH-0104-DAYUAN-FIRST}{前104年·李广利初征大宛失利，汉廷为汗血马再度集兵}账本记录汉与大宛在前104年发生“李广利初征大宛失利，汉廷为汗血马再度集兵”。条目把背景概括为汉廷重视大宛良马以补充骑兵，外交摩擦升级为远距离军事行动，并指出首征暴露河西以西的补给极限，中央投入更多人力物资准备再征。这个节点进入区域社会时，要经由文书、道路、军队、市场、寺院和家庭等不同中介；因此，政治行动的宣布、地方接收和日常生活的改变不能被视作同一时刻完成。

本条依据《史记·大宛列传》；《汉书·张骞李广利传》；《资治通鉴》卷二十，证据提示为“核心事件可靠；军队规模与马种传说有争议”。这些材料可支持年序、官方决定、战事或特定地点的线索，却难以直接统计所有人的财富、迁徙、信仰和动机。更稳妥的读法是区分诏令与实施、军报与人口、行政称谓与自我认同，并让地方材料的缺环限制解释的范围。

由此可见，该事件不是孤立的年表点，而是资源、信息与权力关系重新配置的一环。不同地区的官员、社群和家庭可能以不同方式承担征发、保护、迁移或宗教活动；现存来源只能照亮其中一部分，不能替沉默者制造统一的经验。

\eventanchor{WH-0102-DAYUAN-SECOND}{前102年·李广利再征，大宛杀毋寡并与汉议和}账本记录汉与大宛在前102年发生“李广利再征，大宛杀毋寡并与汉议和”。条目把背景概括为武帝不愿接受首征失败，集中多郡物资并依赖沿途绿洲供给，并指出汉取得马匹和外交威望，却显示西域经营必须依赖脆弱的道路与地方合作。这个节点进入区域社会时，要经由文书、道路、军队、市场、寺院和家庭等不同中介；因此，政治行动的宣布、地方接收和日常生活的改变不能被视作同一时刻完成。

本条依据《史记·大宛列传》；《汉书·张骞李广利传》；《汉书·武帝纪》，证据提示为“年代大体可靠；战果与条约内容有争议”。这些材料可支持年序、官方决定、战事或特定地点的线索，却难以直接统计所有人的财富、迁徙、信仰和动机。更稳妥的读法是区分诏令与实施、军报与人口、行政称谓与自我认同，并让地方材料的缺环限制解释的范围。

由此可见，该事件不是孤立的年表点，而是资源、信息与权力关系重新配置的一环。不同地区的官员、社群和家庭可能以不同方式承担征发、保护、迁移或宗教活动；现存来源只能照亮其中一部分，不能替沉默者制造统一的经验。

\eventanchor{WH-0100-SU-WU}{前100年·苏武奉使匈奴被留，其使节经历成为汉匈关系的象征性叙事}账本记录汉与匈奴在前100年发生“苏武奉使匈奴被留，其使节经历成为汉匈关系的象征性叙事”。条目把背景概括为大宛战后汉匈仍需交涉俘虏和边事，使团处在高度不信任的外交环境，并指出使节扣留反映和战并存的边境政治，苏武故事后来成为官方忠节记忆。这个节点进入区域社会时，要经由文书、道路、军队、市场、寺院和家庭等不同中介；因此，政治行动的宣布、地方接收和日常生活的改变不能被视作同一时刻完成。

本条依据《汉书·苏武传》；《汉书·匈奴传》；《资治通鉴》卷二十，证据提示为“年份和核心事实可靠；细节经后世塑形”。这些材料可支持年序、官方决定、战事或特定地点的线索，却难以直接统计所有人的财富、迁徙、信仰和动机。更稳妥的读法是区分诏令与实施、军报与人口、行政称谓与自我认同，并让地方材料的缺环限制解释的范围。

由此可见，该事件不是孤立的年表点，而是资源、信息与权力关系重新配置的一环。不同地区的官员、社群和家庭可能以不同方式承担征发、保护、迁移或宗教活动；现存来源只能照亮其中一部分，不能替沉默者制造统一的经验。

\eventanchor{WH-0099-LI-LING}{前99年·李陵出塞兵败降匈奴，司马迁为其辩护而受刑}账本记录汉与匈奴在前99年发生“李陵出塞兵败降匈奴，司马迁为其辩护而受刑”。条目把背景概括为武帝持续远征，李陵部在孤军与援兵失约的条件下遭围困，并指出战败引发朝廷责任追究和忠降判断，司马迁受刑影响《史记》的写作处境。这个节点进入区域社会时，要经由文书、道路、军队、市场、寺院和家庭等不同中介；因此，政治行动的宣布、地方接收和日常生活的改变不能被视作同一时刻完成。

本条依据《汉书·李广苏建传》；《史记·太史公自序》；《汉书·司马迁传》，证据提示为“可靠纪年；降因、兵数与司马迁言论细节有争议”。这些材料可支持年序、官方决定、战事或特定地点的线索，却难以直接统计所有人的财富、迁徙、信仰和动机。更稳妥的读法是区分诏令与实施、军报与人口、行政称谓与自我认同，并让地方材料的缺环限制解释的范围。

由此可见，该事件不是孤立的年表点，而是资源、信息与权力关系重新配置的一环。不同地区的官员、社群和家庭可能以不同方式承担征发、保护、迁移或宗教活动；现存来源只能照亮其中一部分，不能替沉默者制造统一的经验。

\eventanchor{WH-0098-COIN-REFORM}{前98年前后·武帝朝持续整顿币制，五铢钱成为主要法定货币}账本记录汉在前98年前后发生“武帝朝持续整顿币制，五铢钱成为主要法定货币”。条目把背景概括为战费和私铸造成币值混乱，中央已通过盐铁和铸币扩张财政控制，并指出五铢钱在帝国范围内取得持久地位，但出土量不能直接等同各地交易规模。这个节点进入区域社会时，要经由文书、道路、军队、市场、寺院和家庭等不同中介；因此，政治行动的宣布、地方接收和日常生活的改变不能被视作同一时刻完成。

本条依据《汉书·食货志》；《史记·平准书》；五铢钱范、钱币窖藏与冶铸遗存，证据提示为“制度趋势可靠；各次改铸与禁令年份有争议”。这些材料可支持年序、官方决定、战事或特定地点的线索，却难以直接统计所有人的财富、迁徙、信仰和动机。更稳妥的读法是区分诏令与实施、军报与人口、行政称谓与自我认同，并让地方材料的缺环限制解释的范围。

由此可见，该事件不是孤立的年表点，而是资源、信息与权力关系重新配置的一环。不同地区的官员、社群和家庭可能以不同方式承担征发、保护、迁移或宗教活动；现存来源只能照亮其中一部分，不能替沉默者制造统一的经验。

\eventanchor{WH-0091-WITCHCRAFT}{前91年·巫蛊案波及太子刘据，长安爆发内战，太子败死}账本记录汉在前91年发生“巫蛊案波及太子刘据，长安爆发内战，太子败死”。条目把背景概括为武帝晚年多疑、江充等人与太子集团冲突，巫蛊指控成为清洗工具，并指出法定继承人死亡，武帝朝战争与财政压力转化为深刻的皇位继承危机。这个节点进入区域社会时，要经由文书、道路、军队、市场、寺院和家庭等不同中介；因此，政治行动的宣布、地方接收和日常生活的改变不能被视作同一时刻完成。

本条依据《汉书·武五子传》；《史记·孝武本纪》；《汉纪》卷二十一，证据提示为“可靠纪年；巫蛊证据与责任分配有争议”。这些材料可支持年序、官方决定、战事或特定地点的线索，却难以直接统计所有人的财富、迁徙、信仰和动机。更稳妥的读法是区分诏令与实施、军报与人口、行政称谓与自我认同，并让地方材料的缺环限制解释的范围。

由此可见，该事件不是孤立的年表点，而是资源、信息与权力关系重新配置的一环。不同地区的官员、社群和家庭可能以不同方式承担征发、保护、迁移或宗教活动；现存来源只能照亮其中一部分，不能替沉默者制造统一的经验。

\eventanchor{WH-0089-LUNTAI}{前89年·武帝下轮台诏，拒绝继续在西域大规模屯田远征}账本记录汉在前89年发生“武帝下轮台诏，拒绝继续在西域大规模屯田远征”。条目把背景概括为连年远征、屯戍和徭役消耗民力，朝廷内部对西域投入的承受力下降，并指出汉廷公开强调休养和慎战，为昭帝、霍光时期的政策收缩提供依据。这个节点进入区域社会时，要经由文书、道路、军队、市场、寺院和家庭等不同中介；因此，政治行动的宣布、地方接收和日常生活的改变不能被视作同一时刻完成。

本条依据《汉书·西域传上》；《汉书·武帝纪》；《资治通鉴》卷二十一，证据提示为“诏书存在可靠；政策转向程度有争议”。这些材料可支持年序、官方决定、战事或特定地点的线索，却难以直接统计所有人的财富、迁徙、信仰和动机。更稳妥的读法是区分诏令与实施、军报与人口、行政称谓与自我认同，并让地方材料的缺环限制解释的范围。

由此可见，该事件不是孤立的年表点，而是资源、信息与权力关系重新配置的一环。不同地区的官员、社群和家庭可能以不同方式承担征发、保护、迁移或宗教活动；现存来源只能照亮其中一部分，不能替沉默者制造统一的经验。

\subsection{边疆、战争与移动}

\eventanchor{WH-0087-WUDI-DEATH}{前87年二月·武帝去世，少帝刘弗陵即位，霍光等受遗诏辅政}账本记录汉在前87年二月发生“武帝去世，少帝刘弗陵即位，霍光等受遗诏辅政”。条目把背景概括为太子刘据已死，新君年幼，武帝须以外戚、将领和官僚共同维持继承，并指出霍光成为辅政核心，政策由扩张转向节制而宫廷权力高度集中。这个节点进入区域社会时，要经由文书、道路、军队、市场、寺院和家庭等不同中介；因此，政治行动的宣布、地方接收和日常生活的改变不能被视作同一时刻完成。

本条依据《汉书·武帝纪》；《汉书·昭帝纪》；《汉书·霍光传》，证据提示为“可靠纪年；遗诏与辅政权力配置有争议”。这些材料可支持年序、官方决定、战事或特定地点的线索，却难以直接统计所有人的财富、迁徙、信仰和动机。更稳妥的读法是区分诏令与实施、军报与人口、行政称谓与自我认同，并让地方材料的缺环限制解释的范围。

由此可见，该事件不是孤立的年表点，而是资源、信息与权力关系重新配置的一环。不同地区的官员、社群和家庭可能以不同方式承担征发、保护、迁移或宗教活动；现存来源只能照亮其中一部分，不能替沉默者制造统一的经验。

\eventanchor{WH-0086-REGENTS}{前86年·霍光、金日磾、上官桀等辅政格局确立}账本记录汉在前86年发生“霍光、金日磾、上官桀等辅政格局确立”。条目把背景概括为幼帝不能亲政，辅政者须控制诏令、禁军和人事以避免重演吕氏危机，并指出霍光的决策优势逐步形成，也埋下上官桀等集团反弹的条件。这个节点进入区域社会时，要经由文书、道路、军队、市场、寺院和家庭等不同中介；因此，政治行动的宣布、地方接收和日常生活的改变不能被视作同一时刻完成。

本条依据《汉书·霍光传》；《汉书·金日磾传》；《汉书·昭帝纪》，证据提示为“相对次序可靠；职权边界有争议”。这些材料可支持年序、官方决定、战事或特定地点的线索，却难以直接统计所有人的财富、迁徙、信仰和动机。更稳妥的读法是区分诏令与实施、军报与人口、行政称谓与自我认同，并让地方材料的缺环限制解释的范围。

由此可见，该事件不是孤立的年表点，而是资源、信息与权力关系重新配置的一环。不同地区的官员、社群和家庭可能以不同方式承担征发、保护、迁移或宗教活动；现存来源只能照亮其中一部分，不能替沉默者制造统一的经验。

\eventanchor{WH-0081-SALT-IRON-DEBATE}{前81年·盐铁会议就专营、边防和财政展开廷议}账本记录汉在前81年发生“盐铁会议就专营、边防和财政展开廷议”。条目把背景概括为武帝财政遗产与边防开支持续，地方儒生和官营体系代表对成本分配存在冲突，并指出专营没有整体废除，会议留下国家经营、民生和边防取舍的系统论辩。这个节点进入区域社会时，要经由文书、道路、军队、市场、寺院和家庭等不同中介；因此，政治行动的宣布、地方接收和日常生活的改变不能被视作同一时刻完成。

本条依据《盐铁论》；《汉书·昭帝纪》；《汉书·食货志》，证据提示为“核心事实可靠；《盐铁论》文本编纂层次有争议”。这些材料可支持年序、官方决定、战事或特定地点的线索，却难以直接统计所有人的财富、迁徙、信仰和动机。更稳妥的读法是区分诏令与实施、军报与人口、行政称谓与自我认同，并让地方材料的缺环限制解释的范围。

由此可见，该事件不是孤立的年表点，而是资源、信息与权力关系重新配置的一环。不同地区的官员、社群和家庭可能以不同方式承担征发、保护、迁移或宗教活动；现存来源只能照亮其中一部分，不能替沉默者制造统一的经验。

\eventanchor{WH-0080-SHANGGUAN-COUP}{前80年·上官桀、燕王刘旦等谋反案败露，霍光巩固辅政权}账本记录汉在前80年发生“上官桀、燕王刘旦等谋反案败露，霍光巩固辅政权”。条目把背景概括为辅政集团内部围绕上官皇后、外戚和皇帝接近权发生竞争，并指出反霍集团被清除，霍光可更直接主导昭帝朝人事和政策。这个节点进入区域社会时，要经由文书、道路、军队、市场、寺院和家庭等不同中介；因此，政治行动的宣布、地方接收和日常生活的改变不能被视作同一时刻完成。

本条依据《汉书·霍光传》；《汉书·燕王刘旦传》；《汉书·昭帝纪》，证据提示为“可靠纪年；密谋情节有争议”。这些材料可支持年序、官方决定、战事或特定地点的线索，却难以直接统计所有人的财富、迁徙、信仰和动机。更稳妥的读法是区分诏令与实施、军报与人口、行政称谓与自我认同，并让地方材料的缺环限制解释的范围。

由此可见，该事件不是孤立的年表点，而是资源、信息与权力关系重新配置的一环。不同地区的官员、社群和家庭可能以不同方式承担征发、保护、迁移或宗教活动；现存来源只能照亮其中一部分，不能替沉默者制造统一的经验。

\eventanchor{WH-0074-ZHAODI-DEATH}{前74年四月·昭帝无嗣去世，霍光迎立昌邑王刘贺}账本记录汉在前74年四月发生“昭帝无嗣去世，霍光迎立昌邑王刘贺”。条目把背景概括为昭帝未留子嗣，霍光须在多支刘氏宗室间迅速确立成年继承人，并指出昌邑王入京使霍光从辅政者转为决定皇位归属的关键人物。这个节点进入区域社会时，要经由文书、道路、军队、市场、寺院和家庭等不同中介；因此，政治行动的宣布、地方接收和日常生活的改变不能被视作同一时刻完成。

本条依据《汉书·昭帝纪》；《汉书·霍光传》；《汉纪》卷二十四，证据提示为“可靠纪年”。这些材料可支持年序、官方决定、战事或特定地点的线索，却难以直接统计所有人的财富、迁徙、信仰和动机。更稳妥的读法是区分诏令与实施、军报与人口、行政称谓与自我认同，并让地方材料的缺环限制解释的范围。

由此可见，该事件不是孤立的年表点，而是资源、信息与权力关系重新配置的一环。不同地区的官员、社群和家庭可能以不同方式承担征发、保护、迁移或宗教活动；现存来源只能照亮其中一部分，不能替沉默者制造统一的经验。

\eventanchor{WH-0074-LIU-HE-DEPOSITION}{前74年六月·霍光以昌邑王刘贺失德为由废之，改立刘病已为宣帝}账本记录汉在前74年六月发生“霍光以昌邑王刘贺失德为由废之，改立刘病已为宣帝”。条目把背景概括为刘贺即位后与霍光集团冲突，霍光掌握太后诏令、朝臣与禁军支持，并指出西汉首次以大臣主导完成皇帝废立，宣帝的即位须与霍光权威共存。这个节点进入区域社会时，要经由文书、道路、军队、市场、寺院和家庭等不同中介；因此，政治行动的宣布、地方接收和日常生活的改变不能被视作同一时刻完成。

本条依据《汉书·霍光传》；《汉书·宣帝纪》；《汉书·昌邑哀王传》，证据提示为“废立事实可靠；失德罪状的叙事有争议”。这些材料可支持年序、官方决定、战事或特定地点的线索，却难以直接统计所有人的财富、迁徙、信仰和动机。更稳妥的读法是区分诏令与实施、军报与人口、行政称谓与自我认同，并让地方材料的缺环限制解释的范围。

由此可见，该事件不是孤立的年表点，而是资源、信息与权力关系重新配置的一环。不同地区的官员、社群和家庭可能以不同方式承担征发、保护、迁移或宗教活动；现存来源只能照亮其中一部分，不能替沉默者制造统一的经验。

\eventanchor{WH-0074-XUANDI-ACCESSION}{前74年六月·刘病已即位，是为宣帝，以民间成长经历进入皇统}账本记录汉在前74年六月发生“刘病已即位，是为宣帝，以民间成长经历进入皇统”。条目把背景概括为霍光需要一位血统明确且政治基础薄弱的宗室继承人，刘病已符合条件，并指出宣帝在霍光监管下继位，后来亲政时形成重视吏治与民生的政治形象。这个节点进入区域社会时，要经由文书、道路、军队、市场、寺院和家庭等不同中介；因此，政治行动的宣布、地方接收和日常生活的改变不能被视作同一时刻完成。

本条依据《汉书·宣帝纪》；《汉书·霍光传》；《汉纪》卷二十四，证据提示为“可靠纪年；早年经历叙事有修饰”。这些材料可支持年序、官方决定、战事或特定地点的线索，却难以直接统计所有人的财富、迁徙、信仰和动机。更稳妥的读法是区分诏令与实施、军报与人口、行政称谓与自我认同，并让地方材料的缺环限制解释的范围。

由此可见，该事件不是孤立的年表点，而是资源、信息与权力关系重新配置的一环。不同地区的官员、社群和家庭可能以不同方式承担征发、保护、迁移或宗教活动；现存来源只能照亮其中一部分，不能替沉默者制造统一的经验。

\eventanchor{WH-0069-XU-EMPEROR-DEATH}{前69年·许皇后遭霍氏家族成员毒害，霍光权势与宫廷婚姻冲突暴露}账本记录汉在前69年发生“许皇后遭霍氏家族成员毒害，霍光权势与宫廷婚姻冲突暴露”。条目把背景概括为霍光家族欲使霍成君为后，宣帝与许氏关系妨碍外戚进一步结合皇权，并指出事件后来成为宣帝清算霍氏的重要理由，外戚与皇帝的信任裂痕加深。这个节点进入区域社会时，要经由文书、道路、军队、市场、寺院和家庭等不同中介；因此，政治行动的宣布、地方接收和日常生活的改变不能被视作同一时刻完成。

本条依据《汉书·外戚传上》；《汉书·霍光传》；《汉书·宣帝纪》，证据提示为“核心事实可靠；具体责任链有争议”。这些材料可支持年序、官方决定、战事或特定地点的线索，却难以直接统计所有人的财富、迁徙、信仰和动机。更稳妥的读法是区分诏令与实施、军报与人口、行政称谓与自我认同，并让地方材料的缺环限制解释的范围。

由此可见，该事件不是孤立的年表点，而是资源、信息与权力关系重新配置的一环。不同地区的官员、社群和家庭可能以不同方式承担征发、保护、迁移或宗教活动；现存来源只能照亮其中一部分，不能替沉默者制造统一的经验。

\eventanchor{WH-0068-HUO-GUANG-DEATH}{前68年·霍光去世，宣帝逐步亲政并削弱霍氏}账本记录汉在前68年发生“霍光去世，宣帝逐步亲政并削弱霍氏”。条目把背景概括为霍光长期控制中枢而其家族继续扩张，宣帝已积累可用的官僚与宗室支持，并指出霍氏后被诛，皇帝重新直接掌控人事与赏罚，辅政政治阶段结束。这个节点进入区域社会时，要经由文书、道路、军队、市场、寺院和家庭等不同中介；因此，政治行动的宣布、地方接收和日常生活的改变不能被视作同一时刻完成。

本条依据《汉书·霍光传》；《汉书·宣帝纪》；《汉纪》卷二十五，证据提示为“可靠纪年；清算先后有争议”。这些材料可支持年序、官方决定、战事或特定地点的线索，却难以直接统计所有人的财富、迁徙、信仰和动机。更稳妥的读法是区分诏令与实施、军报与人口、行政称谓与自我认同，并让地方材料的缺环限制解释的范围。

由此可见，该事件不是孤立的年表点，而是资源、信息与权力关系重新配置的一环。不同地区的官员、社群和家庭可能以不同方式承担征发、保护、迁移或宗教活动；现存来源只能照亮其中一部分，不能替沉默者制造统一的经验。

\eventanchor{WH-0066-JAIL-REFORM}{前66年前后·宣帝重视狱讼，诏令郡国平反冤狱并考核吏治}账本记录汉在前66年前后发生“宣帝重视狱讼，诏令郡国平反冤狱并考核吏治”。条目把背景概括为武帝末和霍光时期的政治案件留下冤狱与吏治压力，宣帝借司法展示亲政，并指出地方官考课与宽刑成为宣帝政治声誉的组成；实际效果仍受地域和档案缺失限制。这个节点进入区域社会时，要经由文书、道路、军队、市场、寺院和家庭等不同中介；因此，政治行动的宣布、地方接收和日常生活的改变不能被视作同一时刻完成。

本条依据《汉书·宣帝纪》；《汉书·循吏传》；《汉书·刑法志》，证据提示为“政策方向可靠；诏令次数与效果有争议”。这些材料可支持年序、官方决定、战事或特定地点的线索，却难以直接统计所有人的财富、迁徙、信仰和动机。更稳妥的读法是区分诏令与实施、军报与人口、行政称谓与自我认同，并让地方材料的缺环限制解释的范围。

由此可见，该事件不是孤立的年表点，而是资源、信息与权力关系重新配置的一环。不同地区的官员、社群和家庭可能以不同方式承担征发、保护、迁移或宗教活动；现存来源只能照亮其中一部分，不能替沉默者制造统一的经验。

\eventanchor{WH-0061-ZHAO-CHONGGUO}{前61年至前60年·赵充国主张屯田、招抚与分化羌部，反对全面进攻}账本记录汉与先零羌在前61年至前60年发生“赵充国主张屯田、招抚与分化羌部，反对全面进攻”。条目把背景概括为羌部起事与边郡军费上升，朝廷在速战与持久经营间争论，并指出屯田和分化策略成为边疆治理案例，但中央对羌部称谓不能替代地方政治的复杂性。这个节点进入区域社会时，要经由文书、道路、军队、市场、寺院和家庭等不同中介；因此，政治行动的宣布、地方接收和日常生活的改变不能被视作同一时刻完成。

本条依据《汉书·赵充国传》；《汉书·宣帝纪》；《汉书·西羌传》，证据提示为“事件年代可靠；政策成效和部众分类有争议”。这些材料可支持年序、官方决定、战事或特定地点的线索，却难以直接统计所有人的财富、迁徙、信仰和动机。更稳妥的读法是区分诏令与实施、军报与人口、行政称谓与自我认同，并让地方材料的缺环限制解释的范围。

由此可见，该事件不是孤立的年表点，而是资源、信息与权力关系重新配置的一环。不同地区的官员、社群和家庭可能以不同方式承担征发、保护、迁移或宗教活动；现存来源只能照亮其中一部分，不能替沉默者制造统一的经验。

\eventanchor{WH-0060-PROTECTOR-WESTERN-REGIONS}{前60年·郑吉受命为西域都护，汉廷设常驻统筹机构}账本记录汉与西域在前60年发生“郑吉受命为西域都护，汉廷设常驻统筹机构”。条目把背景概括为匈奴日逐王降汉后，车师及西域道路的安全条件改善，汉廷需协调驻军和属国，并指出都护成为汉在西域的协调职位，其控制依赖绿洲政权合作和河西补给而非固定疆界。这个节点进入区域社会时，要经由文书、道路、军队、市场、寺院和家庭等不同中介；因此，政治行动的宣布、地方接收和日常生活的改变不能被视作同一时刻完成。

本条依据《汉书·西域传上》；《汉书·宣帝纪》；《资治通鉴》卷二十六，证据提示为“核心事实可靠；都护治所与管辖范围有争议”。这些材料可支持年序、官方决定、战事或特定地点的线索，却难以直接统计所有人的财富、迁徙、信仰和动机。更稳妥的读法是区分诏令与实施、军报与人口、行政称谓与自我认同，并让地方材料的缺环限制解释的范围。

由此可见，该事件不是孤立的年表点，而是资源、信息与权力关系重新配置的一环。不同地区的官员、社群和家庭可能以不同方式承担征发、保护、迁移或宗教活动；现存来源只能照亮其中一部分，不能替沉默者制造统一的经验。

\eventanchor{WH-0058-HUHANYE-CRISIS}{前58年至前56年·匈奴内部分裂，呼韩邪单于寻求与汉结盟}账本记录匈奴与汉在前58年至前56年发生“匈奴内部分裂，呼韩邪单于寻求与汉结盟”。条目把背景概括为匈奴贵族内斗削弱单于权威，呼韩邪面对郅支等对手需要外援，并指出汉廷获得以朝贡、赏赐和军事威慑影响草原政治的机会。这个节点进入区域社会时，要经由文书、道路、军队、市场、寺院和家庭等不同中介；因此，政治行动的宣布、地方接收和日常生活的改变不能被视作同一时刻完成。

本条依据《汉书·匈奴传》；《汉书·萧望之传》；《资治通鉴》卷二十六，证据提示为“相对次序可靠；单于更替日期有争议”。这些材料可支持年序、官方决定、战事或特定地点的线索，却难以直接统计所有人的财富、迁徙、信仰和动机。更稳妥的读法是区分诏令与实施、军报与人口、行政称谓与自我认同，并让地方材料的缺环限制解释的范围。

由此可见，该事件不是孤立的年表点，而是资源、信息与权力关系重新配置的一环。不同地区的官员、社群和家庭可能以不同方式承担征发、保护、迁移或宗教活动；现存来源只能照亮其中一部分，不能替沉默者制造统一的经验。

\eventanchor{WH-0051-HUHANYE-COURT}{前51年·呼韩邪单于入朝称臣，汉廷授予印绶并供给赏赐}账本记录汉与匈奴在前51年发生“呼韩邪单于入朝称臣，汉廷授予印绶并供给赏赐”。条目把背景概括为匈奴内战使呼韩邪需要汉廷支持，汉亦希望稳定北边而避免再度大规模远征，并指出汉匈关系从战争转为以朝贡和资源交换为中心的联盟，草原内部竞争并未结束。这个节点进入区域社会时，要经由文书、道路、军队、市场、寺院和家庭等不同中介；因此，政治行动的宣布、地方接收和日常生活的改变不能被视作同一时刻完成。

本条依据《汉书·匈奴传》；《汉书·宣帝纪》；《汉纪》卷二十七，证据提示为“可靠纪年；礼仪称谓与依附程度有争议”。这些材料可支持年序、官方决定、战事或特定地点的线索，却难以直接统计所有人的财富、迁徙、信仰和动机。更稳妥的读法是区分诏令与实施、军报与人口、行政称谓与自我认同，并让地方材料的缺环限制解释的范围。

由此可见，该事件不是孤立的年表点，而是资源、信息与权力关系重新配置的一环。不同地区的官员、社群和家庭可能以不同方式承担征发、保护、迁移或宗教活动；现存来源只能照亮其中一部分，不能替沉默者制造统一的经验。

\eventanchor{WH-0051-SHIQU-GE}{前51年·石渠阁召开经学会议，讨论五经异同与博士传承}账本记录汉在前51年发生“石渠阁召开经学会议，讨论五经异同与博士传承”。条目把背景概括为博士制度发展出师法家法差异，朝廷需要在选官和礼制中确认可用解释，并指出经学争论被置于宫廷主持下，经典解释与官僚资格的联系加强。这个节点进入区域社会时，要经由文书、道路、军队、市场、寺院和家庭等不同中介；因此，政治行动的宣布、地方接收和日常生活的改变不能被视作同一时刻完成。

本条依据《汉书·宣帝纪》；《汉书·儒林传》；《汉书·艺文志》，证据提示为“会议存在可靠；议题记录不全”。这些材料可支持年序、官方决定、战事或特定地点的线索，却难以直接统计所有人的财富、迁徙、信仰和动机。更稳妥的读法是区分诏令与实施、军报与人口、行政称谓与自我认同，并让地方材料的缺环限制解释的范围。

由此可见，该事件不是孤立的年表点，而是资源、信息与权力关系重新配置的一环。不同地区的官员、社群和家庭可能以不同方式承担征发、保护、迁移或宗教活动；现存来源只能照亮其中一部分，不能替沉默者制造统一的经验。

\eventanchor{WH-0049-XUANDI-DEATH}{前49年·宣帝去世，太子奭即位，是为元帝}账本记录汉在前49年发生“宣帝去世，太子奭即位，是为元帝”。条目把背景概括为宣帝已安排太子继统，但朝廷对刑名、经术和外戚人事存在不同取向，并指出元帝朝政治更倚重儒臣，宣帝时期的吏治与边疆政策面临重新解释。这个节点进入区域社会时，要经由文书、道路、军队、市场、寺院和家庭等不同中介；因此，政治行动的宣布、地方接收和日常生活的改变不能被视作同一时刻完成。

本条依据《汉书·宣帝纪》；《汉书·元帝纪》；《汉纪》卷二十八，证据提示为“可靠纪年”。这些材料可支持年序、官方决定、战事或特定地点的线索，却难以直接统计所有人的财富、迁徙、信仰和动机。更稳妥的读法是区分诏令与实施、军报与人口、行政称谓与自我认同，并让地方材料的缺环限制解释的范围。

由此可见，该事件不是孤立的年表点，而是资源、信息与权力关系重新配置的一环。不同地区的官员、社群和家庭可能以不同方式承担征发、保护、迁移或宗教活动；现存来源只能照亮其中一部分，不能替沉默者制造统一的经验。

\eventanchor{WH-0048-YUANDI-ACCESSION}{前48年·元帝亲政，萧望之等儒臣与宦官、外戚集团竞争}账本记录汉在前48年发生“元帝亲政，萧望之等儒臣与宦官、外戚集团竞争”。条目把背景概括为继位后的用人和诏令渠道决定谁能影响皇帝，宣帝旧臣与新进经生意见不同，并指出宫廷政治重心转向近侍与奏事程序，政策争论常以人物进退的形式呈现。这个节点进入区域社会时，要经由文书、道路、军队、市场、寺院和家庭等不同中介；因此，政治行动的宣布、地方接收和日常生活的改变不能被视作同一时刻完成。

本条依据《汉书·元帝纪》；《汉书·萧望之传》；《汉书·弘恭石显传》，证据提示为“相对次序可靠；派系标签有简化风险”。这些材料可支持年序、官方决定、战事或特定地点的线索，却难以直接统计所有人的财富、迁徙、信仰和动机。更稳妥的读法是区分诏令与实施、军报与人口、行政称谓与自我认同，并让地方材料的缺环限制解释的范围。

由此可见，该事件不是孤立的年表点，而是资源、信息与权力关系重新配置的一环。不同地区的官员、社群和家庭可能以不同方式承担征发、保护、迁移或宗教活动；现存来源只能照亮其中一部分，不能替沉默者制造统一的经验。

\eventanchor{WH-0047-YUANDI-FINANCE}{前47年前后·元帝朝讨论裁减宫室、乐府与部分财政支出}账本记录汉在前47年前后发生“元帝朝讨论裁减宫室、乐府与部分财政支出”。条目把背景概括为儒臣批评武帝以来的奢费和官营负担，灾异解释也推动节用主张，并指出节用成为政治语言，但宫廷、边防和地方财政的实际规模不能由诏令直接判断。这个节点进入区域社会时，要经由文书、道路、军队、市场、寺院和家庭等不同中介；因此，政治行动的宣布、地方接收和日常生活的改变不能被视作同一时刻完成。

本条依据《汉书·元帝纪》；《汉书·贡禹传》；《汉书·食货志》，证据提示为“政策讨论可靠；削减幅度有争议”。这些材料可支持年序、官方决定、战事或特定地点的线索，却难以直接统计所有人的财富、迁徙、信仰和动机。更稳妥的读法是区分诏令与实施、军报与人口、行政称谓与自我认同，并让地方材料的缺环限制解释的范围。

由此可见，该事件不是孤立的年表点，而是资源、信息与权力关系重新配置的一环。不同地区的官员、社群和家庭可能以不同方式承担征发、保护、迁移或宗教活动；现存来源只能照亮其中一部分，不能替沉默者制造统一的经验。

\eventanchor{WH-0044-ZHIZHI}{前44年至前36年·郅支单于西走康居，汉廷持续追踪其西域活动}账本记录汉、匈奴与康居在前44年至前36年发生“郅支单于西走康居，汉廷持续追踪其西域活动”。条目把背景概括为呼韩邪获汉支持后，郅支失去草原竞争优势而向西域和中亚寻求基地，并指出郅支政权使西域都护面临跨越绿洲网络的安全挑战，为陈汤出兵埋下条件。这个节点进入区域社会时，要经由文书、道路、军队、市场、寺院和家庭等不同中介；因此，政治行动的宣布、地方接收和日常生活的改变不能被视作同一时刻完成。

本条依据《汉书·匈奴传》；《汉书·陈汤传》；《汉书·西域传上》，证据提示为“相对次序可靠；路线与康居关系有争议”。这些材料可支持年序、官方决定、战事或特定地点的线索，却难以直接统计所有人的财富、迁徙、信仰和动机。更稳妥的读法是区分诏令与实施、军报与人口、行政称谓与自我认同，并让地方材料的缺环限制解释的范围。

由此可见，该事件不是孤立的年表点，而是资源、信息与权力关系重新配置的一环。不同地区的官员、社群和家庭可能以不同方式承担征发、保护、迁移或宗教活动；现存来源只能照亮其中一部分，不能替沉默者制造统一的经验。

\eventanchor{WH-0036-CHEN-TANG}{前36年·陈汤、甘延寿越权发兵，攻杀郅支单于于康居}账本记录汉、郅支匈奴与康居在前36年发生“陈汤、甘延寿越权发兵，攻杀郅支单于于康居”。条目把背景概括为郅支在康居建立据点并威胁西域交通，前线官员在未获完整诏命下调兵，并指出郅支势力消失，汉廷奖惩越权出兵的同时强化都护对西域军政的象征影响。这个节点进入区域社会时，要经由文书、道路、军队、市场、寺院和家庭等不同中介；因此，政治行动的宣布、地方接收和日常生活的改变不能被视作同一时刻完成。

本条依据《汉书·陈汤传》；《汉书·匈奴传》；《汉书·西域传上》，证据提示为“核心事件可靠；兵力、城址和“明犯强汉者”文字有争议”。这些材料可支持年序、官方决定、战事或特定地点的线索，却难以直接统计所有人的财富、迁徙、信仰和动机。更稳妥的读法是区分诏令与实施、军报与人口、行政称谓与自我认同，并让地方材料的缺环限制解释的范围。

由此可见，该事件不是孤立的年表点，而是资源、信息与权力关系重新配置的一环。不同地区的官员、社群和家庭可能以不同方式承担征发、保护、迁移或宗教活动；现存来源只能照亮其中一部分，不能替沉默者制造统一的经验。

\eventanchor{WH-0033-ZHAOJUN}{前33年·呼韩邪请求汉女为阏氏，王昭君出塞联姻}账本记录汉与匈奴在前33年发生“呼韩邪请求汉女为阏氏，王昭君出塞联姻”。条目把背景概括为呼韩邪朝汉后联盟需以仪式性婚姻巩固，汉廷在宗室婚配之外选取宫女，并指出联姻成为汉匈和平记忆的重要符号，但不能抹去双方各自的政治与资源计算。这个节点进入区域社会时，要经由文书、道路、军队、市场、寺院和家庭等不同中介；因此，政治行动的宣布、地方接收和日常生活的改变不能被视作同一时刻完成。

本条依据《汉书·匈奴传》；《后汉书·南匈奴传》；《资治通鉴》卷二十九，证据提示为“核心事实可靠；昭君身份与后世故事细节有争议”。这些材料可支持年序、官方决定、战事或特定地点的线索，却难以直接统计所有人的财富、迁徙、信仰和动机。更稳妥的读法是区分诏令与实施、军报与人口、行政称谓与自我认同，并让地方材料的缺环限制解释的范围。

由此可见，该事件不是孤立的年表点，而是资源、信息与权力关系重新配置的一环。不同地区的官员、社群和家庭可能以不同方式承担征发、保护、迁移或宗教活动；现存来源只能照亮其中一部分，不能替沉默者制造统一的经验。

\eventanchor{WH-0032-YUANDI-DEATH}{前32年·元帝去世，太子骜即位，是为成帝，王氏外戚地位上升}账本记录汉在前32年发生“元帝去世，太子骜即位，是为成帝，王氏外戚地位上升”。条目把背景概括为元帝后期近侍政治削弱部分儒臣，皇太后王氏家族已拥有显著人脉，并指出成帝朝王凤等相继执政，外戚对中枢人事的影响扩大。这个节点进入区域社会时，要经由文书、道路、军队、市场、寺院和家庭等不同中介；因此，政治行动的宣布、地方接收和日常生活的改变不能被视作同一时刻完成。

本条依据《汉书·元帝纪》；《汉书·成帝纪》；《汉书·外戚传上》，证据提示为“可靠纪年”。这些材料可支持年序、官方决定、战事或特定地点的线索，却难以直接统计所有人的财富、迁徙、信仰和动机。更稳妥的读法是区分诏令与实施、军报与人口、行政称谓与自我认同，并让地方材料的缺环限制解释的范围。

由此可见，该事件不是孤立的年表点，而是资源、信息与权力关系重新配置的一环。不同地区的官员、社群和家庭可能以不同方式承担征发、保护、迁移或宗教活动；现存来源只能照亮其中一部分，不能替沉默者制造统一的经验。

\eventanchor{WH-0028-WANG-FENG}{前28年·王凤长期领尚书事并为大司马，王氏外戚形成中枢集团}账本记录汉在前28年发生“王凤长期领尚书事并为大司马，王氏外戚形成中枢集团”。条目把背景概括为成帝倚重母族且宰相制度与尚书奏事相互交织，王凤积累官僚和宗族资源，并指出王氏由单一后族转为可持续的政治集团，为王莽后来上升提供结构条件。这个节点进入区域社会时，要经由文书、道路、军队、市场、寺院和家庭等不同中介；因此，政治行动的宣布、地方接收和日常生活的改变不能被视作同一时刻完成。

本条依据《汉书·王莽传上》；《汉书·元后传》；《汉书·百官公卿表》，证据提示为“官职与时序可靠；权力机制有争议”。这些材料可支持年序、官方决定、战事或特定地点的线索，却难以直接统计所有人的财富、迁徙、信仰和动机。更稳妥的读法是区分诏令与实施、军报与人口、行政称谓与自我认同，并让地方材料的缺环限制解释的范围。

由此可见，该事件不是孤立的年表点，而是资源、信息与权力关系重新配置的一环。不同地区的官员、社群和家庭可能以不同方式承担征发、保护、迁移或宗教活动；现存来源只能照亮其中一部分，不能替沉默者制造统一的经验。

\eventanchor{WH-0020-CHENGDI-HEIR}{前20年前后·成帝无嗣问题加剧，后宫、外戚与宗室继承安排成为朝廷焦点}账本记录汉在前20年前后发生“成帝无嗣问题加剧，后宫、外戚与宗室继承安排成为朝廷焦点”。条目把背景概括为皇帝长期无成年继承人，王氏、赵氏及诸宗室均可能影响未来皇统，并指出继承的不确定性削弱成帝朝权威，使元帝支系和外戚更深介入储位安排。这个节点进入区域社会时，要经由文书、道路、军队、市场、寺院和家庭等不同中介；因此，政治行动的宣布、地方接收和日常生活的改变不能被视作同一时刻完成。

本条依据《汉书·成帝纪》；《汉书·外戚传下》；《汉书·赵皇后传》，证据提示为“过程性记录可靠；具体因果与医疗叙事有争议”。这些材料可支持年序、官方决定、战事或特定地点的线索，却难以直接统计所有人的财富、迁徙、信仰和动机。更稳妥的读法是区分诏令与实施、军报与人口、行政称谓与自我认同，并让地方材料的缺环限制解释的范围。

由此可见，该事件不是孤立的年表点，而是资源、信息与权力关系重新配置的一环。不同地区的官员、社群和家庭可能以不同方式承担征发、保护、迁移或宗教活动；现存来源只能照亮其中一部分，不能替沉默者制造统一的经验。

\eventanchor{WH-0018-ZHAO-FEIYAN}{前18年·赵飞燕立为皇后，赵氏与王氏围绕后宫资源竞争}账本记录汉在前18年发生“赵飞燕立为皇后，赵氏与王氏围绕后宫资源竞争”。条目把背景概括为成帝宠幸赵氏，后位可改变近侍、封爵和皇嗣保护的资源分配，并指出赵氏获得显要地位但未解决无嗣危机；相关传说不可替代《汉书》纪传证据。这个节点进入区域社会时，要经由文书、道路、军队、市场、寺院和家庭等不同中介；因此，政治行动的宣布、地方接收和日常生活的改变不能被视作同一时刻完成。

本条依据《汉书·成帝纪》；《汉书·外戚传下》；《赵飞燕外传》仅作后出参照，证据提示为“可靠纪年；人物品评多受后世道德化叙事影响”。这些材料可支持年序、官方决定、战事或特定地点的线索，却难以直接统计所有人的财富、迁徙、信仰和动机。更稳妥的读法是区分诏令与实施、军报与人口、行政称谓与自我认同，并让地方材料的缺环限制解释的范围。

由此可见，该事件不是孤立的年表点，而是资源、信息与权力关系重新配置的一环。不同地区的官员、社群和家庭可能以不同方式承担征发、保护、迁移或宗教活动；现存来源只能照亮其中一部分，不能替沉默者制造统一的经验。

\eventanchor{WH-0016-WANG-MANG-MARQUIS}{前16年·王莽受封新都侯，进入王氏外戚的显著位置}账本记录汉在前16年发生“王莽受封新都侯，进入王氏外戚的显著位置”。条目把背景概括为王凤等王氏长辈相继执政，王莽以礼贤、俭约和经学形象积累名望，并指出王莽获得可见的爵位和官僚网络，为日后以“清名”竞争中枢职位奠基。这个节点进入区域社会时，要经由文书、道路、军队、市场、寺院和家庭等不同中介；因此，政治行动的宣布、地方接收和日常生活的改变不能被视作同一时刻完成。

本条依据《汉书·成帝纪》；《汉书·王莽传上》；《汉书·元后传》，证据提示为“可靠纪年；孝行形象有修辞成分”。这些材料可支持年序、官方决定、战事或特定地点的线索，却难以直接统计所有人的财富、迁徙、信仰和动机。更稳妥的读法是区分诏令与实施、军报与人口、行政称谓与自我认同，并让地方材料的缺环限制解释的范围。

由此可见，该事件不是孤立的年表点，而是资源、信息与权力关系重新配置的一环。不同地区的官员、社群和家庭可能以不同方式承担征发、保护、迁移或宗教活动；现存来源只能照亮其中一部分，不能替沉默者制造统一的经验。

\eventanchor{WH-0008-WANG-MANG-REGENT}{前8年·王莽任大司马，开始主持朝政}账本记录汉在前8年发生“王莽任大司马，开始主持朝政”。条目把背景概括为王氏多位大司马相继执政，王莽兼具太后宗族身份和士人声望，并指出王莽可直接参与人事与奏事流程，外戚权力取得新的个人化中心。这个节点进入区域社会时，要经由文书、道路、军队、市场、寺院和家庭等不同中介；因此，政治行动的宣布、地方接收和日常生活的改变不能被视作同一时刻完成。

本条依据《汉书·成帝纪》；《汉书·王莽传上》；《汉书·百官公卿表》，证据提示为“可靠纪年；实际决策范围有争议”。这些材料可支持年序、官方决定、战事或特定地点的线索，却难以直接统计所有人的财富、迁徙、信仰和动机。更稳妥的读法是区分诏令与实施、军报与人口、行政称谓与自我认同，并让地方材料的缺环限制解释的范围。

由此可见，该事件不是孤立的年表点，而是资源、信息与权力关系重新配置的一环。不同地区的官员、社群和家庭可能以不同方式承担征发、保护、迁移或宗教活动；现存来源只能照亮其中一部分，不能替沉默者制造统一的经验。

\eventanchor{WH-0007-AIDI-ACCESSION}{前7年·成帝去世，定陶王刘欣即位，是为哀帝，王莽一度退居}账本记录汉在前7年发生“成帝去世，定陶王刘欣即位，是为哀帝，王莽一度退居”。条目把背景概括为成帝无嗣，朝廷选择定陶王继统；新帝母族与祖母傅氏要求重组权力，并指出王莽失去大司马职位，显示外戚集团并非不可替代，傅氏、丁氏进入中枢。这个节点进入区域社会时，要经由文书、道路、军队、市场、寺院和家庭等不同中介；因此，政治行动的宣布、地方接收和日常生活的改变不能被视作同一时刻完成。

本条依据《汉书·成帝纪》；《汉书·哀帝纪》；《汉书·王莽传上》，证据提示为“可靠纪年”。这些材料可支持年序、官方决定、战事或特定地点的线索，却难以直接统计所有人的财富、迁徙、信仰和动机。更稳妥的读法是区分诏令与实施、军报与人口、行政称谓与自我认同，并让地方材料的缺环限制解释的范围。

由此可见，该事件不是孤立的年表点，而是资源、信息与权力关系重新配置的一环。不同地区的官员、社群和家庭可能以不同方式承担征发、保护、迁移或宗教活动；现存来源只能照亮其中一部分，不能替沉默者制造统一的经验。

\subsection{社会关系与宗教空间}

\eventanchor{WH-0005-DONG-XIAN}{前5年至前1年·哀帝宠任董贤，授大司马并引发官爵安排争议}账本记录汉在前5年至前1年发生“哀帝宠任董贤，授大司马并引发官爵安排争议”。条目把背景概括为哀帝倚重亲近侍从以制衡傅、丁等外戚和旧王氏势力，并指出董贤快速升迁冲击常规官僚序列，哀帝骤死后其集团迅速失去依托。这个节点进入区域社会时，要经由文书、道路、军队、市场、寺院和家庭等不同中介；因此，政治行动的宣布、地方接收和日常生活的改变不能被视作同一时刻完成。

本条依据《汉书·哀帝纪》；《汉书·佞幸传》；《汉书·董贤传》，证据提示为“过程性记录可靠；关系性质的后世标签不宜当事实”。这些材料可支持年序、官方决定、战事或特定地点的线索，却难以直接统计所有人的财富、迁徙、信仰和动机。更稳妥的读法是区分诏令与实施、军报与人口、行政称谓与自我认同，并让地方材料的缺环限制解释的范围。

由此可见，该事件不是孤立的年表点，而是资源、信息与权力关系重新配置的一环。不同地区的官员、社群和家庭可能以不同方式承担征发、保护、迁移或宗教活动；现存来源只能照亮其中一部分，不能替沉默者制造统一的经验。

\eventanchor{WH-0001-AIDI-DEATH}{前1年六月·哀帝无嗣去世，太皇太后王政君召王莽入朝控制局面}账本记录汉在前1年六月发生“哀帝无嗣去世，太皇太后王政君召王莽入朝控制局面”。条目把背景概括为哀帝未定继承人，董贤和傅丁外戚难以形成稳定联盟，太皇太后掌有名义权威，并指出王莽诛逐董贤集团并重掌禁军和政令，西汉继承危机转入摄政阶段。这个节点进入区域社会时，要经由文书、道路、军队、市场、寺院和家庭等不同中介；因此，政治行动的宣布、地方接收和日常生活的改变不能被视作同一时刻完成。

本条依据《汉书·哀帝纪》；《汉书·王莽传上》；《汉书·元后传》，证据提示为“可靠纪年”。这些材料可支持年序、官方决定、战事或特定地点的线索，却难以直接统计所有人的财富、迁徙、信仰和动机。更稳妥的读法是区分诏令与实施、军报与人口、行政称谓与自我认同，并让地方材料的缺环限制解释的范围。

由此可见，该事件不是孤立的年表点，而是资源、信息与权力关系重新配置的一环。不同地区的官员、社群和家庭可能以不同方式承担征发、保护、迁移或宗教活动；现存来源只能照亮其中一部分，不能替沉默者制造统一的经验。

\eventanchor{WH-0001-PINGDI}{元始元年（1年）·中山王刘衎被立为平帝，王莽以大司马辅政}账本记录汉在元始元年（1年）发生“中山王刘衎被立为平帝，王莽以大司马辅政”。条目把背景概括为王莽控制宫廷后需选择年幼且母族势弱的刘氏宗室，以避免新外戚挑战，并指出王莽以辅政者身份控制人事和诏令，开始将摄政包装为奉行周礼的公共职责。这个节点进入区域社会时，要经由文书、道路、军队、市场、寺院和家庭等不同中介；因此，政治行动的宣布、地方接收和日常生活的改变不能被视作同一时刻完成。

本条依据《汉书·平帝纪》；《汉书·王莽传上》；《汉纪》卷三十四，证据提示为“可靠纪年”。这些材料可支持年序、官方决定、战事或特定地点的线索，却难以直接统计所有人的财富、迁徙、信仰和动机。更稳妥的读法是区分诏令与实施、军报与人口、行政称谓与自我认同，并让地方材料的缺环限制解释的范围。

由此可见，该事件不是孤立的年表点，而是资源、信息与权力关系重新配置的一环。不同地区的官员、社群和家庭可能以不同方式承担征发、保护、迁移或宗教活动；现存来源只能照亮其中一部分，不能替沉默者制造统一的经验。

\eventanchor{WH-0002-CENSUS}{元始二年（2年）·《汉书·地理志》保存元始二年户口统计，反映帝国行政申报}账本记录汉在元始二年（2年）发生“《汉书·地理志》保存元始二年户口统计，反映帝国行政申报”。条目把背景概括为摄政政府需掌握郡国户口和赋役基础，郡国上报被编入地理志式统计，并指出该数字成为研究西汉人口的重要基线，但不能视作无漏报的全国实数，简牍仅能校验局部实践。这个节点进入区域社会时，要经由文书、道路、军队、市场、寺院和家庭等不同中介；因此，政治行动的宣布、地方接收和日常生活的改变不能被视作同一时刻完成。

本条依据《汉书·地理志》；尹湾汉墓简牍《集簿》；居延汉简（簿籍与出入文书），证据提示为“统计记录可靠；总数口径与覆盖率有争议”。这些材料可支持年序、官方决定、战事或特定地点的线索，却难以直接统计所有人的财富、迁徙、信仰和动机。更稳妥的读法是区分诏令与实施、军报与人口、行政称谓与自我认同，并让地方材料的缺环限制解释的范围。

由此可见，该事件不是孤立的年表点，而是资源、信息与权力关系重新配置的一环。不同地区的官员、社群和家庭可能以不同方式承担征发、保护、迁移或宗教活动；现存来源只能照亮其中一部分，不能替沉默者制造统一的经验。

\eventanchor{WH-0004-WANG-MANG-DAUGHTER}{元始四年（4年）·王莽女立为皇后，王氏与幼帝的姻亲关系被制度化}账本记录汉在元始四年（4年）发生“王莽女立为皇后，王氏与幼帝的姻亲关系被制度化”。条目把背景概括为王莽已掌握辅政权力，须以皇后婚姻将王氏利益嵌入刘氏皇统，并指出王莽由外戚辅政者变为皇室姻亲，反对者更难以常规外戚名义约束其权力。这个节点进入区域社会时，要经由文书、道路、军队、市场、寺院和家庭等不同中介；因此，政治行动的宣布、地方接收和日常生活的改变不能被视作同一时刻完成。

本条依据《汉书·平帝纪》；《汉书·王莽传上》；《汉书·元后传》，证据提示为“可靠纪年；礼仪安排有争议”。这些材料可支持年序、官方决定、战事或特定地点的线索，却难以直接统计所有人的财富、迁徙、信仰和动机。更稳妥的读法是区分诏令与实施、军报与人口、行政称谓与自我认同，并让地方材料的缺环限制解释的范围。

由此可见，该事件不是孤立的年表点，而是资源、信息与权力关系重新配置的一环。不同地区的官员、社群和家庭可能以不同方式承担征发、保护、迁移或宗教活动；现存来源只能照亮其中一部分，不能替沉默者制造统一的经验。

\eventanchor{WH-0005-PINGDI-DEATH}{元始五年（5年）·平帝去世，王莽立孺子婴为皇太子并居摄}账本记录汉在元始五年（5年）发生“平帝去世，王莽立孺子婴为皇太子并居摄”。条目把背景概括为平帝将近成年，刘氏外戚可能重新进入权力中心；王莽需要延续摄政安排，并指出幼年孺子婴被立而不即帝位，王莽以“居摄”取得超出辅政者的制度位置。这个节点进入区域社会时，要经由文书、道路、军队、市场、寺院和家庭等不同中介；因此，政治行动的宣布、地方接收和日常生活的改变不能被视作同一时刻完成。

本条依据《汉书·平帝纪》；《汉书·王莽传上》；《汉书·外戚传上》，证据提示为“平帝死期可靠；死因与王莽责任有争议”。这些材料可支持年序、官方决定、战事或特定地点的线索，却难以直接统计所有人的财富、迁徙、信仰和动机。更稳妥的读法是区分诏令与实施、军报与人口、行政称谓与自我认同，并让地方材料的缺环限制解释的范围。

由此可见，该事件不是孤立的年表点，而是资源、信息与权力关系重新配置的一环。不同地区的官员、社群和家庭可能以不同方式承担征发、保护、迁移或宗教活动；现存来源只能照亮其中一部分，不能替沉默者制造统一的经验。

\eventanchor{WH-0006-JUSHE}{居摄元年（6年）·王莽以居摄名义颁行诏令、调整官名并处置反对者}账本记录汉在居摄元年（6年）发生“王莽以居摄名义颁行诏令、调整官名并处置反对者”。条目把背景概括为孺子年幼使王莽可将临时辅政延伸为常设权力，符命与周礼话语提供正当化资源，并指出摄政机构开始改造礼制与官制，支持和反对网络均向地方与宗室扩展。这个节点进入区域社会时，要经由文书、道路、军队、市场、寺院和家庭等不同中介；因此，政治行动的宣布、地方接收和日常生活的改变不能被视作同一时刻完成。

本条依据《汉书·王莽传上》；《汉书·平帝纪》；《汉书·百官公卿表》，证据提示为“居摄事实可靠；各项诏令的执行程度有争议”。这些材料可支持年序、官方决定、战事或特定地点的线索，却难以直接统计所有人的财富、迁徙、信仰和动机。更稳妥的读法是区分诏令与实施、军报与人口、行政称谓与自我认同，并让地方材料的缺环限制解释的范围。

由此可见，该事件不是孤立的年表点，而是资源、信息与权力关系重新配置的一环。不同地区的官员、社群和家庭可能以不同方式承担征发、保护、迁移或宗教活动；现存来源只能照亮其中一部分，不能替沉默者制造统一的经验。

\eventanchor{WH-0008-EMPEROR-ACTING}{初始元年（8年）·王莽接受“假皇帝”“摄皇帝”一类符命与劝进，进一步排空孺子婴}账本记录汉在初始元年（8年）发生“王莽接受“假皇帝”“摄皇帝”一类符命与劝进，进一步排空孺子婴”。条目把背景概括为居摄难以长期解释成年政治权力，王莽集团以祥瑞、符命和群臣奏请制造改朝条件，并指出刘氏皇帝仅保留象征位置，禅代所需的礼仪和官僚支持被公开组织。这个节点进入区域社会时，要经由文书、道路、军队、市场、寺院和家庭等不同中介；因此，政治行动的宣布、地方接收和日常生活的改变不能被视作同一时刻完成。

本条依据《汉书·王莽传上》；《汉书·平帝纪》；《汉纪》卷三十五，证据提示为“核心过程可靠；符命真伪与民意程度有争议”。这些材料可支持年序、官方决定、战事或特定地点的线索，却难以直接统计所有人的财富、迁徙、信仰和动机。更稳妥的读法是区分诏令与实施、军报与人口、行政称谓与自我认同，并让地方材料的缺环限制解释的范围。

由此可见，该事件不是孤立的年表点，而是资源、信息与权力关系重新配置的一环。不同地区的官员、社群和家庭可能以不同方式承担征发、保护、迁移或宗教活动；现存来源只能照亮其中一部分，不能替沉默者制造统一的经验。

\eventanchor{WH-0009-XIN-USURPATION}{始建国元年（9年）·王莽受孺子婴禅让，即皇帝位，改国号新}账本记录新在始建国元年（9年）发生“王莽受孺子婴禅让，即皇帝位，改国号新”。条目把背景概括为王莽已控制摄政、婚姻、官僚和符命叙事，刘氏宗室缺乏可协调的反制中心，并指出西汉皇统终结，新朝开始推行改制；禅让叙事与地方接受程度须分开评估。这个节点进入区域社会时，要经由文书、道路、军队、市场、寺院和家庭等不同中介；因此，政治行动的宣布、地方接收和日常生活的改变不能被视作同一时刻完成。

本条依据《汉书·王莽传上》；《汉书·平帝纪》；《汉纪》卷三十五；新莽钱币与官印材料，证据提示为“可靠纪年；禅让自愿性与符命来源有争议”。这些材料可支持年序、官方决定、战事或特定地点的线索，却难以直接统计所有人的财富、迁徙、信仰和动机。更稳妥的读法是区分诏令与实施、军报与人口、行政称谓与自我认同，并让地方材料的缺环限制解释的范围。

由此可见，该事件不是孤立的年表点，而是资源、信息与权力关系重新配置的一环。不同地区的官员、社群和家庭可能以不同方式承担征发、保护、迁移或宗教活动；现存来源只能照亮其中一部分，不能替沉默者制造统一的经验。

\eventanchor{WH-LEDGER-001-1}{前209年七月·“陈胜、吴广在大泽乡起事”的前期动员与资源准备}账本记录秦末诸起事者在前209年七月发生““陈胜、吴广在大泽乡起事”的前期动员与资源准备”。条目把背景概括为当时的权力、资源与信息约束推动相关过程展开，并指出这一过程为“陈胜、吴广在大泽乡起事”提供可观察的条件。这个节点进入区域社会时，要经由文书、道路、军队、市场、寺院和家庭等不同中介；因此，政治行动的宣布、地方接收和日常生活的改变不能被视作同一时刻完成。

本条依据《史记·陈涉世家》；《汉书·陈胜项籍传》；《资治通鉴》卷七，证据提示为“核心事项已有既有账本来源；本分解节点的参与者、执行范围和时间先后仍须逐案核验”。这些材料可支持年序、官方决定、战事或特定地点的线索，却难以直接统计所有人的财富、迁徙、信仰和动机。更稳妥的读法是区分诏令与实施、军报与人口、行政称谓与自我认同，并让地方材料的缺环限制解释的范围。

由此可见，该事件不是孤立的年表点，而是资源、信息与权力关系重新配置的一环。不同地区的官员、社群和家庭可能以不同方式承担征发、保护、迁移或宗教活动；现存来源只能照亮其中一部分，不能替沉默者制造统一的经验。

\eventanchor{WH-LEDGER-001-2}{前209年七月·“陈胜、吴广在大泽乡起事”的命令、联盟或地方响应}账本记录秦末诸起事者在前209年七月发生““陈胜、吴广在大泽乡起事”的命令、联盟或地方响应”。条目把背景概括为当时的权力、资源与信息约束推动相关过程展开，并指出这一过程为“陈胜、吴广在大泽乡起事”提供可观察的条件。这个节点进入区域社会时，要经由文书、道路、军队、市场、寺院和家庭等不同中介；因此，政治行动的宣布、地方接收和日常生活的改变不能被视作同一时刻完成。

本条依据《史记·陈涉世家》；《汉书·陈胜项籍传》；《资治通鉴》卷七，证据提示为“核心事项已有既有账本来源；本分解节点的参与者、执行范围和时间先后仍须逐案核验”。这些材料可支持年序、官方决定、战事或特定地点的线索，却难以直接统计所有人的财富、迁徙、信仰和动机。更稳妥的读法是区分诏令与实施、军报与人口、行政称谓与自我认同，并让地方材料的缺环限制解释的范围。

由此可见，该事件不是孤立的年表点，而是资源、信息与权力关系重新配置的一环。不同地区的官员、社群和家庭可能以不同方式承担征发、保护、迁移或宗教活动；现存来源只能照亮其中一部分，不能替沉默者制造统一的经验。

\eventanchor{WH-LEDGER-001-3}{前209年七月·陈胜、吴广在大泽乡起事}账本记录秦末诸起事者在前209年七月发生“陈胜、吴广在大泽乡起事”。条目把背景概括为“陈胜、吴广在大泽乡起事”发生的政治与区域条件共同作用，并指出该节点改变后续资源、联盟或制度处置的条件。这个节点进入区域社会时，要经由文书、道路、军队、市场、寺院和家庭等不同中介；因此，政治行动的宣布、地方接收和日常生活的改变不能被视作同一时刻完成。

本条依据《史记·陈涉世家》；《汉书·陈胜项籍传》；《资治通鉴》卷七，证据提示为“核心事项已有既有账本来源；本分解节点的参与者、执行范围和时间先后仍须逐案核验”。这些材料可支持年序、官方决定、战事或特定地点的线索，却难以直接统计所有人的财富、迁徙、信仰和动机。更稳妥的读法是区分诏令与实施、军报与人口、行政称谓与自我认同，并让地方材料的缺环限制解释的范围。

由此可见，该事件不是孤立的年表点，而是资源、信息与权力关系重新配置的一环。不同地区的官员、社群和家庭可能以不同方式承担征发、保护、迁移或宗教活动；现存来源只能照亮其中一部分，不能替沉默者制造统一的经验。

\eventanchor{WH-LEDGER-001-4}{前209年七月·“陈胜、吴广在大泽乡起事”的执行中的空间与人员处置}账本记录秦末诸起事者在前209年七月发生““陈胜、吴广在大泽乡起事”的执行中的空间与人员处置”。条目把背景概括为“陈胜、吴广在大泽乡起事”发生的政治与区域条件共同作用，并指出该节点改变后续资源、联盟或制度处置的条件。这个节点进入区域社会时，要经由文书、道路、军队、市场、寺院和家庭等不同中介；因此，政治行动的宣布、地方接收和日常生活的改变不能被视作同一时刻完成。

本条依据《史记·陈涉世家》；《汉书·陈胜项籍传》；《资治通鉴》卷七，证据提示为“核心事项已有既有账本来源；本分解节点的参与者、执行范围和时间先后仍须逐案核验”。这些材料可支持年序、官方决定、战事或特定地点的线索，却难以直接统计所有人的财富、迁徙、信仰和动机。更稳妥的读法是区分诏令与实施、军报与人口、行政称谓与自我认同，并让地方材料的缺环限制解释的范围。

由此可见，该事件不是孤立的年表点，而是资源、信息与权力关系重新配置的一环。不同地区的官员、社群和家庭可能以不同方式承担征发、保护、迁移或宗教活动；现存来源只能照亮其中一部分，不能替沉默者制造统一的经验。

\eventanchor{WH-LEDGER-001-5}{前209年七月·“陈胜、吴广在大泽乡起事”的后续制度或军政调整}账本记录秦末诸起事者在前209年七月发生““陈胜、吴广在大泽乡起事”的后续制度或军政调整”。条目把背景概括为核心节点发生后仍须处理其遗留问题，并指出后续影响须在不同地区和材料类型中分别核验。这个节点进入区域社会时，要经由文书、道路、军队、市场、寺院和家庭等不同中介；因此，政治行动的宣布、地方接收和日常生活的改变不能被视作同一时刻完成。

本条依据《史记·陈涉世家》；《汉书·陈胜项籍传》；《资治通鉴》卷七，证据提示为“核心事项已有既有账本来源；本分解节点的参与者、执行范围和时间先后仍须逐案核验”。这些材料可支持年序、官方决定、战事或特定地点的线索，却难以直接统计所有人的财富、迁徙、信仰和动机。更稳妥的读法是区分诏令与实施、军报与人口、行政称谓与自我认同，并让地方材料的缺环限制解释的范围。

由此可见，该事件不是孤立的年表点，而是资源、信息与权力关系重新配置的一环。不同地区的官员、社群和家庭可能以不同方式承担征发、保护、迁移或宗教活动；现存来源只能照亮其中一部分，不能替沉默者制造统一的经验。

\eventanchor{WH-LEDGER-001-6}{前209年七月·“陈胜、吴广在大泽乡起事”的异源材料与影响范围的复核}账本记录秦末诸起事者在前209年七月发生““陈胜、吴广在大泽乡起事”的异源材料与影响范围的复核”。条目把背景概括为核心节点发生后仍须处理其遗留问题，并指出后续影响须在不同地区和材料类型中分别核验。这个节点进入区域社会时，要经由文书、道路、军队、市场、寺院和家庭等不同中介；因此，政治行动的宣布、地方接收和日常生活的改变不能被视作同一时刻完成。

本条依据《史记·陈涉世家》；《汉书·陈胜项籍传》；《资治通鉴》卷七，证据提示为“核心事项已有既有账本来源；本分解节点的参与者、执行范围和时间先后仍须逐案核验”。这些材料可支持年序、官方决定、战事或特定地点的线索，却难以直接统计所有人的财富、迁徙、信仰和动机。更稳妥的读法是区分诏令与实施、军报与人口、行政称谓与自我认同，并让地方材料的缺环限制解释的范围。

由此可见，该事件不是孤立的年表点，而是资源、信息与权力关系重新配置的一环。不同地区的官员、社群和家庭可能以不同方式承担征发、保护、迁移或宗教活动；现存来源只能照亮其中一部分，不能替沉默者制造统一的经验。

\eventanchor{WH-LEDGER-002-1}{前208年·“项梁拥立楚怀王后，重建楚国名号”的前期动员与资源准备}账本记录楚地反秦集团在前208年发生““项梁拥立楚怀王后，重建楚国名号”的前期动员与资源准备”。条目把背景概括为当时的权力、资源与信息约束推动相关过程展开，并指出这一过程为“项梁拥立楚怀王后，重建楚国名号”提供可观察的条件。这个节点进入区域社会时，要经由文书、道路、军队、市场、寺院和家庭等不同中介；因此，政治行动的宣布、地方接收和日常生活的改变不能被视作同一时刻完成。

本条依据《史记·项羽本纪》；《史记·陈涉世家》；《汉书·陈胜项籍传》，证据提示为“核心事项已有既有账本来源；本分解节点的参与者、执行范围和时间先后仍须逐案核验”。这些材料可支持年序、官方决定、战事或特定地点的线索，却难以直接统计所有人的财富、迁徙、信仰和动机。更稳妥的读法是区分诏令与实施、军报与人口、行政称谓与自我认同，并让地方材料的缺环限制解释的范围。

由此可见，该事件不是孤立的年表点，而是资源、信息与权力关系重新配置的一环。不同地区的官员、社群和家庭可能以不同方式承担征发、保护、迁移或宗教活动；现存来源只能照亮其中一部分，不能替沉默者制造统一的经验。

\eventanchor{WH-LEDGER-002-2}{前208年·“项梁拥立楚怀王后，重建楚国名号”的命令、联盟或地方响应}账本记录楚地反秦集团在前208年发生““项梁拥立楚怀王后，重建楚国名号”的命令、联盟或地方响应”。条目把背景概括为当时的权力、资源与信息约束推动相关过程展开，并指出这一过程为“项梁拥立楚怀王后，重建楚国名号”提供可观察的条件。这个节点进入区域社会时，要经由文书、道路、军队、市场、寺院和家庭等不同中介；因此，政治行动的宣布、地方接收和日常生活的改变不能被视作同一时刻完成。

本条依据《史记·项羽本纪》；《史记·陈涉世家》；《汉书·陈胜项籍传》，证据提示为“核心事项已有既有账本来源；本分解节点的参与者、执行范围和时间先后仍须逐案核验”。这些材料可支持年序、官方决定、战事或特定地点的线索，却难以直接统计所有人的财富、迁徙、信仰和动机。更稳妥的读法是区分诏令与实施、军报与人口、行政称谓与自我认同，并让地方材料的缺环限制解释的范围。

由此可见，该事件不是孤立的年表点，而是资源、信息与权力关系重新配置的一环。不同地区的官员、社群和家庭可能以不同方式承担征发、保护、迁移或宗教活动；现存来源只能照亮其中一部分，不能替沉默者制造统一的经验。

\eventanchor{WH-LEDGER-002-3}{前208年·项梁拥立楚怀王后，重建楚国名号}账本记录楚地反秦集团在前208年发生“项梁拥立楚怀王后，重建楚国名号”。条目把背景概括为“项梁拥立楚怀王后，重建楚国名号”发生的政治与区域条件共同作用，并指出该节点改变后续资源、联盟或制度处置的条件。这个节点进入区域社会时，要经由文书、道路、军队、市场、寺院和家庭等不同中介；因此，政治行动的宣布、地方接收和日常生活的改变不能被视作同一时刻完成。

本条依据《史记·项羽本纪》；《史记·陈涉世家》；《汉书·陈胜项籍传》，证据提示为“核心事项已有既有账本来源；本分解节点的参与者、执行范围和时间先后仍须逐案核验”。这些材料可支持年序、官方决定、战事或特定地点的线索，却难以直接统计所有人的财富、迁徙、信仰和动机。更稳妥的读法是区分诏令与实施、军报与人口、行政称谓与自我认同，并让地方材料的缺环限制解释的范围。

由此可见，该事件不是孤立的年表点，而是资源、信息与权力关系重新配置的一环。不同地区的官员、社群和家庭可能以不同方式承担征发、保护、迁移或宗教活动；现存来源只能照亮其中一部分，不能替沉默者制造统一的经验。

\eventanchor{WH-LEDGER-002-4}{前208年·“项梁拥立楚怀王后，重建楚国名号”的执行中的空间与人员处置}账本记录楚地反秦集团在前208年发生““项梁拥立楚怀王后，重建楚国名号”的执行中的空间与人员处置”。条目把背景概括为“项梁拥立楚怀王后，重建楚国名号”发生的政治与区域条件共同作用，并指出该节点改变后续资源、联盟或制度处置的条件。这个节点进入区域社会时，要经由文书、道路、军队、市场、寺院和家庭等不同中介；因此，政治行动的宣布、地方接收和日常生活的改变不能被视作同一时刻完成。

本条依据《史记·项羽本纪》；《史记·陈涉世家》；《汉书·陈胜项籍传》，证据提示为“核心事项已有既有账本来源；本分解节点的参与者、执行范围和时间先后仍须逐案核验”。这些材料可支持年序、官方决定、战事或特定地点的线索，却难以直接统计所有人的财富、迁徙、信仰和动机。更稳妥的读法是区分诏令与实施、军报与人口、行政称谓与自我认同，并让地方材料的缺环限制解释的范围。

由此可见，该事件不是孤立的年表点，而是资源、信息与权力关系重新配置的一环。不同地区的官员、社群和家庭可能以不同方式承担征发、保护、迁移或宗教活动；现存来源只能照亮其中一部分，不能替沉默者制造统一的经验。

\eventanchor{WH-LEDGER-002-5}{前208年·“项梁拥立楚怀王后，重建楚国名号”的后续制度或军政调整}账本记录楚地反秦集团在前208年发生““项梁拥立楚怀王后，重建楚国名号”的后续制度或军政调整”。条目把背景概括为核心节点发生后仍须处理其遗留问题，并指出后续影响须在不同地区和材料类型中分别核验。这个节点进入区域社会时，要经由文书、道路、军队、市场、寺院和家庭等不同中介；因此，政治行动的宣布、地方接收和日常生活的改变不能被视作同一时刻完成。

本条依据《史记·项羽本纪》；《史记·陈涉世家》；《汉书·陈胜项籍传》，证据提示为“核心事项已有既有账本来源；本分解节点的参与者、执行范围和时间先后仍须逐案核验”。这些材料可支持年序、官方决定、战事或特定地点的线索，却难以直接统计所有人的财富、迁徙、信仰和动机。更稳妥的读法是区分诏令与实施、军报与人口、行政称谓与自我认同，并让地方材料的缺环限制解释的范围。

由此可见，该事件不是孤立的年表点，而是资源、信息与权力关系重新配置的一环。不同地区的官员、社群和家庭可能以不同方式承担征发、保护、迁移或宗教活动；现存来源只能照亮其中一部分，不能替沉默者制造统一的经验。

\eventanchor{WH-LEDGER-002-6}{前208年·“项梁拥立楚怀王后，重建楚国名号”的异源材料与影响范围的复核}账本记录楚地反秦集团在前208年发生““项梁拥立楚怀王后，重建楚国名号”的异源材料与影响范围的复核”。条目把背景概括为核心节点发生后仍须处理其遗留问题，并指出后续影响须在不同地区和材料类型中分别核验。这个节点进入区域社会时，要经由文书、道路、军队、市场、寺院和家庭等不同中介；因此，政治行动的宣布、地方接收和日常生活的改变不能被视作同一时刻完成。

本条依据《史记·项羽本纪》；《史记·陈涉世家》；《汉书·陈胜项籍传》，证据提示为“核心事项已有既有账本来源；本分解节点的参与者、执行范围和时间先后仍须逐案核验”。这些材料可支持年序、官方决定、战事或特定地点的线索，却难以直接统计所有人的财富、迁徙、信仰和动机。更稳妥的读法是区分诏令与实施、军报与人口、行政称谓与自我认同，并让地方材料的缺环限制解释的范围。

由此可见，该事件不是孤立的年表点，而是资源、信息与权力关系重新配置的一环。不同地区的官员、社群和家庭可能以不同方式承担征发、保护、迁移或宗教活动；现存来源只能照亮其中一部分，不能替沉默者制造统一的经验。

\eventanchor{WH-LEDGER-003-1}{前207年·“巨鹿之战中项羽军击破秦军主力”的前期动员与资源准备}账本记录楚军与秦军在前207年发生““巨鹿之战中项羽军击破秦军主力”的前期动员与资源准备”。条目把背景概括为当时的权力、资源与信息约束推动相关过程展开，并指出这一过程为“巨鹿之战中项羽军击破秦军主力”提供可观察的条件。这个节点进入区域社会时，要经由文书、道路、军队、市场、寺院和家庭等不同中介；因此，政治行动的宣布、地方接收和日常生活的改变不能被视作同一时刻完成。

本条依据《史记·项羽本纪》；《史记·秦始皇本纪》；《汉书·陈胜项籍传》；《资治通鉴》卷八，证据提示为“核心事项已有既有账本来源；本分解节点的参与者、执行范围和时间先后仍须逐案核验”。这些材料可支持年序、官方决定、战事或特定地点的线索，却难以直接统计所有人的财富、迁徙、信仰和动机。更稳妥的读法是区分诏令与实施、军报与人口、行政称谓与自我认同，并让地方材料的缺环限制解释的范围。

由此可见，该事件不是孤立的年表点，而是资源、信息与权力关系重新配置的一环。不同地区的官员、社群和家庭可能以不同方式承担征发、保护、迁移或宗教活动；现存来源只能照亮其中一部分，不能替沉默者制造统一的经验。

\eventanchor{WH-LEDGER-003-2}{前207年·“巨鹿之战中项羽军击破秦军主力”的命令、联盟或地方响应}账本记录楚军与秦军在前207年发生““巨鹿之战中项羽军击破秦军主力”的命令、联盟或地方响应”。条目把背景概括为当时的权力、资源与信息约束推动相关过程展开，并指出这一过程为“巨鹿之战中项羽军击破秦军主力”提供可观察的条件。这个节点进入区域社会时，要经由文书、道路、军队、市场、寺院和家庭等不同中介；因此，政治行动的宣布、地方接收和日常生活的改变不能被视作同一时刻完成。

本条依据《史记·项羽本纪》；《史记·秦始皇本纪》；《汉书·陈胜项籍传》；《资治通鉴》卷八，证据提示为“核心事项已有既有账本来源；本分解节点的参与者、执行范围和时间先后仍须逐案核验”。这些材料可支持年序、官方决定、战事或特定地点的线索，却难以直接统计所有人的财富、迁徙、信仰和动机。更稳妥的读法是区分诏令与实施、军报与人口、行政称谓与自我认同，并让地方材料的缺环限制解释的范围。

由此可见，该事件不是孤立的年表点，而是资源、信息与权力关系重新配置的一环。不同地区的官员、社群和家庭可能以不同方式承担征发、保护、迁移或宗教活动；现存来源只能照亮其中一部分，不能替沉默者制造统一的经验。

\eventanchor{WH-LEDGER-003-3}{前207年·巨鹿之战中项羽军击破秦军主力}账本记录楚军与秦军在前207年发生“巨鹿之战中项羽军击破秦军主力”。条目把背景概括为“巨鹿之战中项羽军击破秦军主力”发生的政治与区域条件共同作用，并指出该节点改变后续资源、联盟或制度处置的条件。这个节点进入区域社会时，要经由文书、道路、军队、市场、寺院和家庭等不同中介；因此，政治行动的宣布、地方接收和日常生活的改变不能被视作同一时刻完成。

本条依据《史记·项羽本纪》；《史记·秦始皇本纪》；《汉书·陈胜项籍传》；《资治通鉴》卷八，证据提示为“核心事项已有既有账本来源；本分解节点的参与者、执行范围和时间先后仍须逐案核验”。这些材料可支持年序、官方决定、战事或特定地点的线索，却难以直接统计所有人的财富、迁徙、信仰和动机。更稳妥的读法是区分诏令与实施、军报与人口、行政称谓与自我认同，并让地方材料的缺环限制解释的范围。

由此可见，该事件不是孤立的年表点，而是资源、信息与权力关系重新配置的一环。不同地区的官员、社群和家庭可能以不同方式承担征发、保护、迁移或宗教活动；现存来源只能照亮其中一部分，不能替沉默者制造统一的经验。

\eventanchor{WH-LEDGER-003-4}{前207年·“巨鹿之战中项羽军击破秦军主力”的执行中的空间与人员处置}账本记录楚军与秦军在前207年发生““巨鹿之战中项羽军击破秦军主力”的执行中的空间与人员处置”。条目把背景概括为“巨鹿之战中项羽军击破秦军主力”发生的政治与区域条件共同作用，并指出该节点改变后续资源、联盟或制度处置的条件。这个节点进入区域社会时，要经由文书、道路、军队、市场、寺院和家庭等不同中介；因此，政治行动的宣布、地方接收和日常生活的改变不能被视作同一时刻完成。

本条依据《史记·项羽本纪》；《史记·秦始皇本纪》；《汉书·陈胜项籍传》；《资治通鉴》卷八，证据提示为“核心事项已有既有账本来源；本分解节点的参与者、执行范围和时间先后仍须逐案核验”。这些材料可支持年序、官方决定、战事或特定地点的线索，却难以直接统计所有人的财富、迁徙、信仰和动机。更稳妥的读法是区分诏令与实施、军报与人口、行政称谓与自我认同，并让地方材料的缺环限制解释的范围。

由此可见，该事件不是孤立的年表点，而是资源、信息与权力关系重新配置的一环。不同地区的官员、社群和家庭可能以不同方式承担征发、保护、迁移或宗教活动；现存来源只能照亮其中一部分，不能替沉默者制造统一的经验。

\eventanchor{WH-LEDGER-003-5}{前207年·“巨鹿之战中项羽军击破秦军主力”的后续制度或军政调整}账本记录楚军与秦军在前207年发生““巨鹿之战中项羽军击破秦军主力”的后续制度或军政调整”。条目把背景概括为核心节点发生后仍须处理其遗留问题，并指出后续影响须在不同地区和材料类型中分别核验。这个节点进入区域社会时，要经由文书、道路、军队、市场、寺院和家庭等不同中介；因此，政治行动的宣布、地方接收和日常生活的改变不能被视作同一时刻完成。

本条依据《史记·项羽本纪》；《史记·秦始皇本纪》；《汉书·陈胜项籍传》；《资治通鉴》卷八，证据提示为“核心事项已有既有账本来源；本分解节点的参与者、执行范围和时间先后仍须逐案核验”。这些材料可支持年序、官方决定、战事或特定地点的线索，却难以直接统计所有人的财富、迁徙、信仰和动机。更稳妥的读法是区分诏令与实施、军报与人口、行政称谓与自我认同，并让地方材料的缺环限制解释的范围。

由此可见，该事件不是孤立的年表点，而是资源、信息与权力关系重新配置的一环。不同地区的官员、社群和家庭可能以不同方式承担征发、保护、迁移或宗教活动；现存来源只能照亮其中一部分，不能替沉默者制造统一的经验。

\eventanchor{WH-LEDGER-003-6}{前207年·“巨鹿之战中项羽军击破秦军主力”的异源材料与影响范围的复核}账本记录楚军与秦军在前207年发生““巨鹿之战中项羽军击破秦军主力”的异源材料与影响范围的复核”。条目把背景概括为核心节点发生后仍须处理其遗留问题，并指出后续影响须在不同地区和材料类型中分别核验。这个节点进入区域社会时，要经由文书、道路、军队、市场、寺院和家庭等不同中介；因此，政治行动的宣布、地方接收和日常生活的改变不能被视作同一时刻完成。

本条依据《史记·项羽本纪》；《史记·秦始皇本纪》；《汉书·陈胜项籍传》；《资治通鉴》卷八，证据提示为“核心事项已有既有账本来源；本分解节点的参与者、执行范围和时间先后仍须逐案核验”。这些材料可支持年序、官方决定、战事或特定地点的线索，却难以直接统计所有人的财富、迁徙、信仰和动机。更稳妥的读法是区分诏令与实施、军报与人口、行政称谓与自我认同，并让地方材料的缺环限制解释的范围。

由此可见，该事件不是孤立的年表点，而是资源、信息与权力关系重新配置的一环。不同地区的官员、社群和家庭可能以不同方式承担征发、保护、迁移或宗教活动；现存来源只能照亮其中一部分，不能替沉默者制造统一的经验。

\eventanchor{WH-LEDGER-004-1}{前207年冬·“章邯向项羽降，秦军降卒后遭坑杀的叙事引发争议”的前期动员与资源准备}账本记录秦与楚在前207年冬发生““章邯向项羽降，秦军降卒后遭坑杀的叙事引发争议”的前期动员与资源准备”。条目把背景概括为当时的权力、资源与信息约束推动相关过程展开，并指出这一过程为“章邯向项羽降，秦军降卒后遭坑杀的叙事引发争议”提供可观察的条件。这个节点进入区域社会时，要经由文书、道路、军队、市场、寺院和家庭等不同中介；因此，政治行动的宣布、地方接收和日常生活的改变不能被视作同一时刻完成。

本条依据《史记·项羽本纪》；《汉书·陈胜项籍传》；《资治通鉴》卷八，证据提示为“核心事项已有既有账本来源；本分解节点的参与者、执行范围和时间先后仍须逐案核验”。这些材料可支持年序、官方决定、战事或特定地点的线索，却难以直接统计所有人的财富、迁徙、信仰和动机。更稳妥的读法是区分诏令与实施、军报与人口、行政称谓与自我认同，并让地方材料的缺环限制解释的范围。

由此可见，该事件不是孤立的年表点，而是资源、信息与权力关系重新配置的一环。不同地区的官员、社群和家庭可能以不同方式承担征发、保护、迁移或宗教活动；现存来源只能照亮其中一部分，不能替沉默者制造统一的经验。
